% Options for packages loaded elsewhere
\PassOptionsToPackage{unicode}{hyperref}
\PassOptionsToPackage{hyphens}{url}
\documentclass[
]{article}
\usepackage{xcolor}
\usepackage{amsmath,amssymb}
\setcounter{secnumdepth}{-\maxdimen} % remove section numbering
\usepackage{iftex}
\ifPDFTeX
  \usepackage[T1]{fontenc}
  \usepackage[utf8]{inputenc}
  \usepackage{textcomp} % provide euro and other symbols
\else % if luatex or xetex
  \usepackage{unicode-math} % this also loads fontspec
  \defaultfontfeatures{Scale=MatchLowercase}
  \defaultfontfeatures[\rmfamily]{Ligatures=TeX,Scale=1}
\fi
\usepackage{lmodern}
\ifPDFTeX\else
  % xetex/luatex font selection
\fi
% Use upquote if available, for straight quotes in verbatim environments
\IfFileExists{upquote.sty}{\usepackage{upquote}}{}
\IfFileExists{microtype.sty}{% use microtype if available
  \usepackage[]{microtype}
  \UseMicrotypeSet[protrusion]{basicmath} % disable protrusion for tt fonts
}{}
\makeatletter
\@ifundefined{KOMAClassName}{% if non-KOMA class
  \IfFileExists{parskip.sty}{%
    \usepackage{parskip}
  }{% else
    \setlength{\parindent}{0pt}
    \setlength{\parskip}{6pt plus 2pt minus 1pt}}
}{% if KOMA class
  \KOMAoptions{parskip=half}}
\makeatother
\usepackage{color}
\usepackage{fancyvrb}
\newcommand{\VerbBar}{|}
\newcommand{\VERB}{\Verb[commandchars=\\\{\}]}
\DefineVerbatimEnvironment{Highlighting}{Verbatim}{commandchars=\\\{\}}
% Add ',fontsize=\small' for more characters per line
\newenvironment{Shaded}{}{}
\newcommand{\AlertTok}[1]{\textcolor[rgb]{1.00,0.00,0.00}{\textbf{#1}}}
\newcommand{\AnnotationTok}[1]{\textcolor[rgb]{0.38,0.63,0.69}{\textbf{\textit{#1}}}}
\newcommand{\AttributeTok}[1]{\textcolor[rgb]{0.49,0.56,0.16}{#1}}
\newcommand{\BaseNTok}[1]{\textcolor[rgb]{0.25,0.63,0.44}{#1}}
\newcommand{\BuiltInTok}[1]{\textcolor[rgb]{0.00,0.50,0.00}{#1}}
\newcommand{\CharTok}[1]{\textcolor[rgb]{0.25,0.44,0.63}{#1}}
\newcommand{\CommentTok}[1]{\textcolor[rgb]{0.38,0.63,0.69}{\textit{#1}}}
\newcommand{\CommentVarTok}[1]{\textcolor[rgb]{0.38,0.63,0.69}{\textbf{\textit{#1}}}}
\newcommand{\ConstantTok}[1]{\textcolor[rgb]{0.53,0.00,0.00}{#1}}
\newcommand{\ControlFlowTok}[1]{\textcolor[rgb]{0.00,0.44,0.13}{\textbf{#1}}}
\newcommand{\DataTypeTok}[1]{\textcolor[rgb]{0.56,0.13,0.00}{#1}}
\newcommand{\DecValTok}[1]{\textcolor[rgb]{0.25,0.63,0.44}{#1}}
\newcommand{\DocumentationTok}[1]{\textcolor[rgb]{0.73,0.13,0.13}{\textit{#1}}}
\newcommand{\ErrorTok}[1]{\textcolor[rgb]{1.00,0.00,0.00}{\textbf{#1}}}
\newcommand{\ExtensionTok}[1]{#1}
\newcommand{\FloatTok}[1]{\textcolor[rgb]{0.25,0.63,0.44}{#1}}
\newcommand{\FunctionTok}[1]{\textcolor[rgb]{0.02,0.16,0.49}{#1}}
\newcommand{\ImportTok}[1]{\textcolor[rgb]{0.00,0.50,0.00}{\textbf{#1}}}
\newcommand{\InformationTok}[1]{\textcolor[rgb]{0.38,0.63,0.69}{\textbf{\textit{#1}}}}
\newcommand{\KeywordTok}[1]{\textcolor[rgb]{0.00,0.44,0.13}{\textbf{#1}}}
\newcommand{\NormalTok}[1]{#1}
\newcommand{\OperatorTok}[1]{\textcolor[rgb]{0.40,0.40,0.40}{#1}}
\newcommand{\OtherTok}[1]{\textcolor[rgb]{0.00,0.44,0.13}{#1}}
\newcommand{\PreprocessorTok}[1]{\textcolor[rgb]{0.74,0.48,0.00}{#1}}
\newcommand{\RegionMarkerTok}[1]{#1}
\newcommand{\SpecialCharTok}[1]{\textcolor[rgb]{0.25,0.44,0.63}{#1}}
\newcommand{\SpecialStringTok}[1]{\textcolor[rgb]{0.73,0.40,0.53}{#1}}
\newcommand{\StringTok}[1]{\textcolor[rgb]{0.25,0.44,0.63}{#1}}
\newcommand{\VariableTok}[1]{\textcolor[rgb]{0.10,0.09,0.49}{#1}}
\newcommand{\VerbatimStringTok}[1]{\textcolor[rgb]{0.25,0.44,0.63}{#1}}
\newcommand{\WarningTok}[1]{\textcolor[rgb]{0.38,0.63,0.69}{\textbf{\textit{#1}}}}
\usepackage{longtable,booktabs,array}
\newcounter{none} % for unnumbered tables
\usepackage{calc} % for calculating minipage widths
% Correct order of tables after \paragraph or \subparagraph
\usepackage{etoolbox}
\makeatletter
\patchcmd\longtable{\par}{\if@noskipsec\mbox{}\fi\par}{}{}
\makeatother
% Allow footnotes in longtable head/foot
\IfFileExists{footnotehyper.sty}{\usepackage{footnotehyper}}{\usepackage{footnote}}
\makesavenoteenv{longtable}
\usepackage{graphicx}
\makeatletter
\newsavebox\pandoc@box
\newcommand*\pandocbounded[1]{% scales image to fit in text height/width
  \sbox\pandoc@box{#1}%
  \Gscale@div\@tempa{\textheight}{\dimexpr\ht\pandoc@box+\dp\pandoc@box\relax}%
  \Gscale@div\@tempb{\linewidth}{\wd\pandoc@box}%
  \ifdim\@tempb\p@<\@tempa\p@\let\@tempa\@tempb\fi% select the smaller of both
  \ifdim\@tempa\p@<\p@\scalebox{\@tempa}{\usebox\pandoc@box}%
  \else\usebox{\pandoc@box}%
  \fi%
}
% Set default figure placement to htbp
\def\fps@figure{htbp}
\makeatother
\setlength{\emergencystretch}{3em} % prevent overfull lines
\providecommand{\tightlist}{%
  \setlength{\itemsep}{0pt}\setlength{\parskip}{0pt}}
\usepackage{bookmark}
\IfFileExists{xurl.sty}{\usepackage{xurl}}{} % add URL line breaks if available
\urlstyle{same}
\hypersetup{
  hidelinks,
  pdfcreator={LaTeX via pandoc}}

\author{}
\date{}

\begin{document}

\section{Runtime Security for Agentic
Systems}\label{runtime-security-for-agentic-systems}

\subsection{A Practical Two-Plane Approach for OpenClaw-Class
Agents}\label{a-practical-two-plane-approach-for-openclaw-class-agents}

Authors: Frank Lyonnet, Antoine Clerget\\
Status: arXiv draft (not submitted)\\
Version: 2026-03-02

\subsection{Abstract}\label{abstract}

Tool-using AI agents expose a runtime security gap: pre-install scanning
and configuration hardening are necessary, but they cannot fully
characterize behavior that emerges during execution (e.g., indirect
prompt injection and tool poisoning). We present a two-plane security
model that correlates reasoning-plane observations (independently read
session transcripts) with system-plane telemetry to detect
intent-behavior divergence.

The contribution is architectural: process-attributed telemetry,
task-scoped intent declarations, and a confidence-stratified signal
hierarchy form a portable control loop across MCP-compatible stacks. In
the reference implementation, EDAMAME accepts agent-tagged
behavioral-model slices from multiple reasoning-plane producers,
including OpenClaw and Cursor, merges them into one observer-owned
correlation window, and preserves contributor attribution. To
demonstrate feasibility and relevance, the live benchmark in this draft
uses an OpenClaw-derived scenario corpus and the explicit-slice producer
path (\texttt{upsert\_behavioral\_model}) to exercise EDAMAME's
downstream correlation logic, while the same EDAMAME control loop also
supports raw-session ingest
(\texttt{upsert\_behavioral\_model\_from\_raw\_sessions}) for Cursor on
a developer workstation.

On the current canonical BadAgentUse artifact lineage (50 live scenarios
across six benchmark categories, plus two integration-validated
categories; 273 valid runs across 12 seeds), precision is 0.862 (CI95:
{[}0.815, 0.899{]}) and recall is 0.987 (CI95: {[}0.962, 0.996{]}).
Median detection latency is 45.8 s (p95: 53.6 s): fast enough to support
many operator responses, but still a meaningful pre-alert window and
therefore the primary limitation of a reactive design. These aggregate
metrics are cross-checked by layered model-to-verdict validation of the
core engine and a focused live validation ladder spanning clean,
unexpected-egress, \texttt{/tmp} lineage, and designed-miss behavior. We
explicitly mark designed blind spots: living-off-the-land behavior under
allowed process/destination pairs and timing-window evasion before
telemetry arrives. This frames the model as a practical correlation
architecture with an explicit latency trade-off versus
interception-first systems such as AgentSentinel, while still providing
a portable and trace-auditable runtime detection path.

\subsection{Keywords}\label{keywords}

agent security, runtime monitoring, prompt injection, tool poisoning,
Model Context Protocol (MCP), intent-behavior correlation, telemetry,
reproducibility

\subsection{1. Introduction}\label{introduction}

Agent systems now perform high-impact operations through tools and MCP
integrations: network access, file mutation, credential-adjacent
actions, and persistent memory updates. Risk therefore depends on what
the agent actually does at runtime, not only on what is statically
configured before runtime.

A recent publicly discussed incident involving Meta security leadership
and an OpenClaw deployment that reportedly executed unintended
mailbox-deletion actions illustrates the practical problem: loss of
control can happen during normal runtime, even in technically
sophisticated environments. Our model targets that
runtime-loss-of-control phase by providing independent host-side
observation and divergence alerts. It does not replace authorization
design or pre-action approval gates; if broad destructive permissions
are already granted, prevention still depends on those controls.

We characterize the problem in terms of three \textbf{control strata} to
avoid terminology collision with our model's two planes:

\begin{enumerate}
\def\labelenumi{\arabic{enumi}.}
\tightlist
\item
  \textbf{S1: Supply-chain controls} (e.g., signed artifacts,
  provenance/SBOM checks).
\item
  \textbf{S2: Configuration controls} (e.g., localhost MCP bind + PSK
  auth, role and tool scoping, sandboxing, policy hardening).
\item
  \textbf{S3: Runtime behavior controls} (e.g., a declared read-only
  task followed by undeclared egress, detected via intent-behavior
  consistency checks).
\end{enumerate}

The residual runtime-only failures that remain after S1/S2 controls are
what we call the \textbf{S3 gap}.

Together, S1/S2 can stop many failures. For example, an internet-exposed
OpenClaw gateway (SecurityScorecard STRIKE, 2026) is often prevented by
correct bind/auth configuration, which is fundamentally a
configuration-control problem. However, S1/S2 do not fully address
runtime divergence: in CVE-2026-25253, a browser token exfiltration path
can execute during normal user interaction even when baseline setup
appears acceptable. Similar runtime-only gaps arise in coding agents
operating over production repositories, where unauthorized access to
git-managed secrets, SSH keys, or CI tokens can produce outsized impact.

This paper addresses the S3 gap with a two-plane model: compare what the
reasoning plane declares with what the system plane observes on the host
(defined in Section 3). The approach highlights a structural opportunity
for security tooling vendors (here, EDAMAME Posture) and AI platform
providers (here, OpenClaw and Cursor) to collaborate through MCP: the
security tool provides system-plane ground truth that no AI framework
includes natively, while the AI platform provides the runtime
observation surface (session transcripts, memory, tool invocation
history) that no external monitor can access alone. Neither plane is
sufficient in isolation; the two-plane correlation is the security
primitive. By design, this first-generation rule set has a blind spot
for living-off-the-land behavior when both process and destination are
already authorized; we surface that limitation explicitly in Sections 7
and 9.

\subsubsection{1.1 Contributions (Novelty)}\label{contributions-novelty}

We separate contributions into a model part and an
implementation/evaluation part.

\textbf{Model contributions (Part I):}

\begin{itemize}
\tightlist
\item
  \textbf{Two-plane runtime security model.} A portable control loop for
  correlating observed agent behavior (reasoning plane, via session
  transcript observation) with independent host telemetry (system
  plane).
\item
  \textbf{Model-level requirements.} Process-attributed telemetry and a
  signal hierarchy are presented as requirements of the model, not as
  product-specific implementation details.
\item
  \textbf{MCP-native framing and portability.} We argue that MCP is not
  only an integration transport but a key enabler of framework-agnostic
  runtime verification.
\end{itemize}

\textbf{Implementation and evaluation contributions (Part II):}

\begin{itemize}
\tightlist
\item
  \textbf{OpenClaw + Cursor-capable EDAMAME instantiation.} A concrete
  deployment path over local, PSK-authenticated MCP with
  rollback-capable action handling and merged multi-agent
  reasoning-plane inputs. The quantitative benchmark in this draft uses
  OpenClaw-derived scenarios with explicit prebuilt-window
  behavioral-model slices published into EDAMAME, while Cursor is
  integrated through the raw-session observer contract rather than
  treated as a separate detection engine.
\item
  \textbf{Trace-backed benchmark protocol.} Every verdict is tied to
  on-disk telemetry artifacts for auditability.
\item
  \textbf{Empirical performance characterization.} Precision/recall,
  latency, and failure modes are measured on a live 50-scenario suite,
  including explicit designed-miss scenarios, and cross-checked with
  layered model-to-verdict validation.
\end{itemize}

\subsection{Part I --- Two-Plane Model}\label{part-i-two-plane-model}

Part I is organized as follows: Section 2 positions this work in prior
runtime security literature; Section 3 defines the two-plane model;
Section 4 states threat assumptions and boundaries; and Section 5 gives
concise properties of the decision rules as a bridge to implementation.
Readers primarily interested in model mechanics and applications may
skip Section 2 and proceed directly to Section 3.

\subsection{2. Background and Related
Work}\label{background-and-related-work}

\subsubsection{2.1 Attack Surface}\label{attack-surface}

Indirect prompt injection blurs data and instructions in LLM-integrated
applications, enabling remote compromise via retrieved content (Greshake
et al., 2023). Benchmarks such as InjecAgent show tool-integrated agents
remain vulnerable to indirect prompt injection (Zhan et al., 2024).

MCP expands the attack surface by standardizing tool descriptors and
invocations. Attacks embedded in tool metadata (tool poisoning,
descriptor mutation) achieve attack success rates up to 72.8\% on
capable models (MCPTox; Luo et al., 2025). Mitigations include
descriptor signing, semantic vetting (Jamshidi et al., 2025), and
multi-stage neural detection pipelines achieving 96\% accuracy on
adversarial prompts (MCP-Guard; Chen et al., 2025). MCP's own
specification documents authorization and security best practices,
including confused deputy risks in proxying designs (MCP Spec, 2025).
MCP Security Bench provides a 2,000-instance, 12-attack-type evaluation
suite for systematically measuring tool-level defenses (Wang et al.,
2025).

Goal drift -- an agent's tendency to deviate from original objectives
over extended execution -- is now measurable in controlled evaluations.
Zhou et al. (2025) show that all evaluated models exhibit some drift
under autonomous operation, even when strong scaffolding substantially
reduces it.

\subsubsection{2.2 Runtime Defense
Landscape}\label{runtime-defense-landscape}

The 2025-2026 research cycle has produced a rich landscape of runtime
defense systems. We organize them by their primary signal source.

\textbf{Reasoning-plane defenses} monitor the agent's internal state
(prompts, chain of thought, tool call sequences) to detect compromise:

\begin{itemize}
\tightlist
\item
  AgentArmor (Gao et al., 2025) treats runtime traces as program
  objects, constructs CFG/DFG/PDG views, and applies type-checking to
  trust-boundary flows. It reports 97\% defense success (3\% attack
  success) on AgentDojo with modest utility impact.
\item
  LlamaFirewall (Meta, 2025) provides three production guardrails:
  PromptGuard 2 (BERT-based jailbreak detection), Agent Alignment Checks
  (chain-of-thought auditing for prompt injection and goal
  misalignment), and CodeShield (online static analysis). It is deployed
  in production at Meta.
\item
  AegisLLM (Li et al., 2025) uses cooperative multi-agent defense where
  orchestrator, deflector, responder, and evaluator agents collaborate
  for safe outputs, reporting a 51\% improvement on jailbreak benchmarks
  without retraining.
\item
  A2AS (Wallarm et al., 2025) proposes the BASIC security model:
  behavior certificates, authenticated prompts, security boundaries,
  in-context defenses, and codified policies. The initiative has broad
  industry sponsorship and is emerging as a framework-level governance
  model.
\end{itemize}

\textbf{Hybrid defenses} combine reasoning-state signals and runtime
action signals:

\begin{itemize}
\tightlist
\item
  MI9 (Singh et al., 2025) is the first integrated runtime governance
  framework, combining goal-conditioned drift detection, FSM-based
  conformance engines, continuous authorization monitoring, and
  graduated containment. Reports 99.81\% detection on synthetic
  scenarios.
\item
  AgentSentinel (Hu et al., 2025) intercepts all sensitive operations
  within agent services, halting execution until a security audit
  correlating task context with system traces completes. Achieves 79.6\%
  defense success across 60 attack scenarios in 6 categories.
\item
  Pro2Guard (Zhang et al., 2025) learns Discrete-Time Markov Chains
  (DTMCs) from execution traces and proactively predicts unsafe states
  before they occur, with PAC-style correctness bounds. Reported defense
  success is 97.4\% (2.6\% unsafe outcomes).
\end{itemize}

\textbf{Formal verification approaches} provide provable safety
properties:

\begin{itemize}
\tightlist
\item
  ShieldAgent (Xu et al., 2025) constructs safety policies as verifiable
  rule circuits and generates shielding plans with executable code for
  formal verification. Achieves 90.1\% recall on safety violations.
\item
  Agent-C (Kamath et al., 2025) enforces temporal safety constraints via
  SMT solving and constrained generation, achieving 100\% conformance on
  leading-model evaluations.
\item
  VET (Li et al., 2025) establishes verifiable execution traces via
  trusted hardware, cryptographic proofs, and notarized TLS transcripts
  with practical overhead under 3x for API calls.
\item
  AgentGuard (Koohestani, 2025) provides runtime verification via
  probabilistic model checking and Dynamic Probabilistic Assurance for
  quantitative agent properties.
\end{itemize}

\textbf{Chain-of-thought monitoring research} reveals important
limitations:

\begin{itemize}
\tightlist
\item
  Models can generate benign-sounding CoT explanations that mask
  misaligned intentions, causing monitors to fail (Thought Crime; Rogers
  et al., 2025).
\item
  Agent awareness of monitoring is the single strongest evasion
  advantage -- detection rates drop sharply when agents know they are
  monitored (Baker et al., 2025).
\item
  Hybrid protocols scoring both reasoning and actions consistently
  surpass either alone, with 2x higher detection for subtle deception
  scenarios (CoT Red-Handed; Goldowsky-Dill et al., 2025).
\end{itemize}

\subsubsection{2.3 Positioning}\label{positioning}

Table 2 summarizes the landscape. Our system provides MCP-native,
multiplatform process-attributed telemetry integrated with intent
correlation. The closest comparator is AgentSentinel, which also
correlates task context with OS-level traces via eBPF/LSM. Key
differentiators:

\begin{itemize}
\tightlist
\item
  \textbf{Multiplatform}: our observer (EDAMAME Posture) works across
  Linux, macOS, and Windows (Section 4.3), while AgentSentinel requires
  Linux eBPF/LSM. eBPF enhances attribution resolution on Linux but is
  not required.
\item
  \textbf{MCP-native}: telemetry is exposed as standard MCP tools over
  Streamable HTTP with PSK authentication, usable by any MCP-compatible
  agent without custom integration.
\item
  \textbf{Correlation vs.~interception}: AgentSentinel halts execution
  at sensitive operations for synchronous audit; our system primarily
  correlates observed behavior after it appears in telemetry, with an
  optional proactive layer for MCP-gated actions. This trades immediate
  blocking for lower runtime friction and broader platform
  compatibility.
\item
  \textbf{Latency trade-off (explicit).} Our measured 45.8 s median
  detection delay (p95: 53.6 s) is a practical limitation of reactive
  observation; AgentSentinel reduces this gap for intercepted operations
  but at the cost of inline interception and Linux/eBPF coupling.
\end{itemize}

Its primary limitations relative to the field are partial
reasoning-plane coverage (transcript-only, no CoT introspection; Section
9) and absence of formal safety proofs.

\begingroup\small
\begin{center}
\setlength{\tabcolsep}{4pt}
\renewcommand{\arraystretch}{1.04}
\begin{tabular}{@{}>{\raggedright\arraybackslash}p{0.18\linewidth}>{\raggedright\arraybackslash}p{0.14\linewidth}>{\raggedright\arraybackslash}p{0.08\linewidth}>{\raggedright\arraybackslash}p{0.18\linewidth}>{\raggedright\arraybackslash}p{0.30\linewidth}@{}}
\toprule
System & Plane & MCP & Platforms & Eval. Scale \\
\midrule
AgentArmor & Reasoning & No & Any & AgentDojo (97\%) \\
LlamaFirewall & Reasoning & No & Any & Meta internal \\
AegisLLM & Reasoning & No & Any & Jailbreak bench. \\
A2AS BASIC & Reasoning & Yes & Any & Framework-level \\
MI9 & Hybrid & No & Any & Synthetic (99.81\%) \\
AgentSentinel & Hybrid & No & Linux (eBPF) & 60 scen., 6 cat. \\
Pro2Guard & Hybrid & No & Any & Household + AV \\
ShieldAgent & Formal & No & Any & Web environments \\
Agent-C & Formal & No & Any & Multiple bench. \\
VET & Trace & No & Any & Black-box API \\
MCP-Guard & Metadata & Yes & Any & 70,448 samples \\
\textbf{This work} & \textbf{Hybrid} & \textbf{Yes} & \textbf{Lin/Mac/Win} & \textbf{50 scen., 8 cat.} \\
\bottomrule
\end{tabular}
\end{center}
\endgroup

\begingroup\footnotesize\noindent A2AS BASIC is conceptual at the
MCP-native layer. \texttt{This work} is hybrid because reasoning-plane
sessions and system-plane telemetry are observed independently and
merged only inside the observer-owned runtime.

\par\endgroup

Formal properties reported by prior systems: - A2AS (behavioral
certificates) - MI9 (FSM conformance) - Pro2Guard (PAC-correctness) -
ShieldAgent (rule circuits) - Agent-C (SMT temporal constraints) - VET
(cryptographic proof) - MCP-Guard (neural detector, reported 96\%)

Deployed systems in this survey scope: - LlamaFirewall (Meta production)
- A2AS (emerging standard) - This work (Lima VM) - All others are
currently research-only artifacts.

\begin{figure}
\centering
\pandocbounded{\includegraphics[keepaspectratio,alt={Agent runtime security timeline}]{figures/fig6_timeline.pdf}}
\caption{Agent runtime security timeline}
\end{figure}

No prior published defense, to our knowledge, exposes continual,
process-attributed OS-level telemetry (eBPF sessions, parent process
lineage, open file paths) as MCP tools and then demonstrates
intent-telemetry correlation through that same interface.\footnote{Survey
  scope (as of February 2026): arXiv, ACM DL, IEEE Xplore, and USENIX
  proceedings through February 2026 using keywords including ``agent
  runtime monitoring'', ``MCP security'', ``tool-use telemetry'', and
  ``eBPF agent'', plus manual review of MCP specifications and ecosystem
  tooling.} Close neighbors (AgentSentinel, LlamaFirewall, MCP-Guard)
intercept or audit tool calls but do not expose OS-level telemetry as a
first-class MCP self-inspection surface. The trade-off is that
system-plane monitoring is inherently reactive (Section 3.6) and
reasoning-plane coverage is transcript-only, without CoT introspection
(Section 9).

\subsubsection{2.4 Reflection as Relevant Prior
Art}\label{reflection-as-relevant-prior-art}

The model is also related to reflection-oriented programming, where a
program observes and reasons about its own behavior through a meta-level
interface. Sobel and Friedman (1996) frame reflection as separating
base-level execution from meta-level supervision. Our two-plane model
adopts a similar separation: the reasoning plane executes and declares
intent; the system plane independently observes effects and can
challenge that declaration. We do not claim formal equivalence to
classic reflection systems, but the analogy is useful for conceptual
grounding.

\subsection{3. Two-Plane Architecture (Reasoning Plane and System
Plane)}\label{two-plane-architecture-reasoning-plane-and-system-plane}

\begin{itemize}
\tightlist
\item
  \textbf{Reasoning plane:} what the agent declares it is doing and why.
\item
  \textbf{System plane:} what the host actually does.
\end{itemize}

At a high level, the control loop is:

declare intent -\textgreater{}\\
derive expected effects -\textgreater{}\\
observe host behavior -\textgreater{}\\
correlate -\textgreater{}\\
compute deterministic evidence -\textgreater{}\\
apply LLM adjudication with critical-signal guardrails -\textgreater{}\\
emit final divergence verdict

\begin{figure}
\centering
\pandocbounded{\includegraphics[keepaspectratio,alt={Two-plane architecture with reasoning-plane declarations and system-plane runtime loops.}]{figures/fig1_architecture.pdf}}
\caption{Two-plane architecture with reasoning-plane declarations and
system-plane runtime loops.}
\end{figure}

The architecture maps reasoning-plane intent declarations to
system-plane telemetry correlation, and emits divergence when observed
behavior is not explained by declared intent. In the deployed system,
the behavioral model consumed by the divergence engine can be a merged
window assembled from multiple agent-tagged contributors (for example
OpenClaw plus Cursor) rather than a single producer. Producer modes are
additive: OpenClaw commonly uses the prebuilt-window publication path,
while workstation integrations such as Cursor can forward raw reasoning
sessions so EDAMAME generates the contributor slice internally before
merge.

The operational cadence is shared-timer-driven: extrapolator and
divergence verdict-reader jobs run on offset timers, while EDAMAME
evaluates three adjacent runtime loops internally. The timer structure
is important to portability: it avoids dependence on framework-specific
synchronous hooks and keeps behavior consistent across deployments. We
considered two alternatives -- interception-first synchronous gating
(lower delay but higher runtime friction and tighter platform coupling)
and a single merged polling loop (simpler but weaker separation between
model publication and verdict consumption). The shared-timer split gave
the best operator clarity while preserving framework-agnostic
deployment.

The paper's model claim remains two-plane correlation. The three-loop
split is an implementation decomposition of the system plane in the
reference instantiation:

\begin{enumerate}
\def\labelenumi{\arabic{enumi}.}
\tightlist
\item
  \textbf{Divergence / intent loop.} Correlates the merged behavioral
  model against raw process-attributed telemetry and emits
  intent-relative evidence when observations are not explained by
  declared behavior.
\item
  \textbf{Vulnerability / safety-floor loop.} Evaluates
  model-independent guardrail findings and runtime incident checks (for
  example \texttt{/tmp} lineage, credential-file-plus-egress, exposed
  listeners, and other vulnerability detectors), including cases where
  the behavioral model is missing or stale.
\item
  \textbf{Advisor / remediation loop.} Groups posture, breach, LAN, and
  high-confidence telemetry findings into operator-facing todos and
  reversible actions.
\end{enumerate}

This separation is deliberate: telemetry confidence labels such as
\texttt{anomalous} or \texttt{blacklisted} are useful operational
signals, but they are not treated as intent-relative divergence evidence
by themselves. Scenarios are judged from engine verdicts read via
\texttt{get\_divergence\_verdict}; the final operator verdict may be
LLM-adjudicated, but deterministic safety-floor and vulnerability-loop
findings remain non-downgradable.

\begin{figure}
\centering
\pandocbounded{\includegraphics[keepaspectratio,alt={Shared-timer cadence for extrapolator publication, internal runtime loops, and verdict consumption.}]{figures/fig8_cron_sequence.pdf}}
\caption{Shared-timer cadence for extrapolator publication, internal
runtime loops, and verdict consumption.}
\end{figure}

\begin{figure}
\centering
\pandocbounded{\includegraphics[keepaspectratio,alt={Reference runtime decomposition of the three adjacent loops inside the system plane.}]{figures/fig9_divergence_engine.pdf}}
\caption{Reference runtime decomposition of the three adjacent loops
inside the system plane.}
\end{figure}

\subsubsection{3.1 Reasoning-Plane Inputs}\label{reasoning-plane-inputs}

The reasoning plane provides structured declarations:

\begin{itemize}
\tightlist
\item
  intent statement (task-scoped)
\item
  tool plan and rationale
\item
  predicted side effects
\item
  task-conditioned allowlists (processes/domains/paths)
\end{itemize}

\textbf{Session-transcript observation (reference implementation, Part
II).} This detail matters because it changes the trust model of the
reasoning plane. In the reference implementation, the reasoning plane is
not self-declared. For OpenClaw, a separate monitor session reads other
sessions' tool calls and reasoning via the platform's
\texttt{sessions\_list} and \texttt{sessions\_history} APIs. For Cursor,
a parallel publisher emits the same behavioral-window schema with its
own \texttt{agent\_type} and \texttt{agent\_instance\_id}. EDAMAME
stores these slices, preserves contributor attribution, and merges them
before correlation. This means an agent cannot retroactively edit or
suppress evidence of its tool calls when the transcript or observer
surface is platform-controlled. The reasoning plane therefore shifts
from trust-dependent self-declaration to an independently verifiable
record, though transcript integrity still depends on the platform (see
Section 9, Limitations).

\subsubsection{3.2 System-Plane
Observations}\label{system-plane-observations}

The system plane consumes independent host telemetry:

\begin{itemize}
\tightlist
\item
  network sessions with \textbf{Layer-7 process enrichment} (binary,
  parent lineage, open files, command line, resource usage)
\item
  telemetry confidence labels (\texttt{exceptions}, \texttt{anomalous},
  \texttt{blacklisted})
\item
  host LAN visibility and local open-port observations
\item
  posture score deltas and threat-model outputs
\end{itemize}

These telemetry labels are organized as a strict confidence hierarchy:

\begin{Shaded}
\begin{Highlighting}[]
\NormalTok{sessions {-}\textgreater{} exceptions {-}\textgreater{} anomalous {-}\textgreater{} blacklisted}
\end{Highlighting}
\end{Shaded}

\begin{figure}
\centering
\pandocbounded{\includegraphics[keepaspectratio,alt={Signal hierarchy with escalating detection confidence}]{figures/fig7_signal_hierarchy.pdf}}
\caption{Signal hierarchy with escalating detection confidence}
\end{figure}

The hierarchy is part of the model as a system-plane
evidence-prioritization mechanism, not merely an EDAMAME implementation
detail: it separates ``did anything happen?'' from ``how strongly should
we trust that this is malicious?'' and allows conservative policy tuning
at lower tiers. In the reference implementation, however, those labels
do not replace intent correlation. The divergence / intent loop reasons
primarily over raw sessions plus Layer-7, LAN, and open-port context;
the vulnerability / safety-floor loop and the advisor / remediation loop
consume \texttt{anomalous} and \texttt{blacklisted} labels as
high-confidence operational signals.

\subsubsection{3.3 Why Process Attribution Is
Model-Critical}\label{why-process-attribution-is-model-critical}

The two-plane model requires process attribution to remain meaningful:

\begin{itemize}
\tightlist
\item
  without process attribution, an allowlist collapses to
  destination-only checks, so benign and malicious programs contacting
  the same endpoint become indistinguishable
\item
  parent-process lineage is required to detect spawn-chain manipulation,
  where an approved binary is launched through an untrusted wrapper
  chain
\item
  open-file context is required for credential-adjacent detection,
  because outbound traffic alone cannot show whether sensitive files
  were read just before egress
\end{itemize}

This requirement is independent of implementation. Any observer that
cannot attribute connections to concrete processes cannot fully realize
the model.

\begin{figure}
\centering
\pandocbounded{\includegraphics[keepaspectratio,alt={Layer-7 Session Enrichment Schema.}]{figures/fig3_l7_schema.pdf}}
\caption{Layer-7 Session Enrichment Schema.}
\end{figure}

Each session record is enriched with process identity, parent lineage,
security signals (including \texttt{open\_files} and
\texttt{spawned\_from\_tmp}), and resource metrics. In this schema,
safety-floor rules apply across the full evidence bundle, not only to a
single branch.

\subsubsection{3.4 Why MCP Matters to the
Model}\label{why-mcp-matters-to-the-model}

MCP is not only a transport convenience. It makes the model
operationally portable:

\begin{enumerate}
\def\labelenumi{\arabic{enumi}.}
\tightlist
\item
  \textbf{Structured semantics.} Typed responses reduce parsing
  ambiguity and let correlation logic consume machine-readable evidence
  directly.
\item
  \textbf{Transactional actions.} Execute/rollback patterns make
  pre-action and post-action checks composable.
\item
  \textbf{Framework portability.} MCP-facing tools can be consumed by
  OpenClaw, Cursor, other coding agents, and other MCP-compatible
  clients with minimal glue code. In the deployed contract, each
  producer tags its slice with agent identity and EDAMAME performs the
  merge centrally.
\end{enumerate}

As positioned in Section 2.3, exposing continual process-attributed host
telemetry as first-class MCP tools is central to this instantiation.

\subsubsection{3.5 Multiplatform Requirement and
Instantiation}\label{multiplatform-requirement-and-instantiation}

Multiplatform operation is not cosmetic. Agent workflows already span
Linux servers, macOS developer hosts, and Windows enterprise endpoints.
This is especially true for coding agents operating across local repos,
CI runners, and production deployment environments.

Accordingly, the model requires attribution beyond Linux-only eBPF
paths. Our current implementation uses:

\begin{itemize}
\tightlist
\item
  Linux eBPF where available
\item
  native macOS/Windows process APIs
\item
  a socket-matching fallback when kernel telemetry is unavailable
\end{itemize}

\begin{figure}
\centering
\pandocbounded{\includegraphics[keepaspectratio,alt={Multiplatform Process Attribution: Linux uses eBPF kprobes for real-time attribution, macOS uses proc\_pidinfo, and Windows uses NtQuerySystemInformation. All platforms share a universal fallback via netstat2 socket matching with port cache.}]{figures/fig2_multiplatform.pdf}}
\caption{Multiplatform Process Attribution: Linux uses eBPF kprobes for
real-time attribution, macOS uses proc\_pidinfo, and Windows uses
NtQuerySystemInformation. All platforms share a universal fallback via
netstat2 socket matching with port cache.}
\end{figure}

\subsubsection{3.6 Observation Delay Is a Constraint, Not a
Contribution}\label{observation-delay-is-a-constraint-not-a-contribution}

Telemetry is delayed relative to action execution. We therefore use a
settling schedule (e.g., 15 s -\textgreater{} 30 s -\textgreater{} 60 s)
and stop when snapshots stabilize. This is an operational accommodation,
not a conceptual novelty claim.

\subsubsection{3.7 Two Example Rules (Simple and
Realistic)}\label{two-example-rules-simple-and-realistic}

We present two example rules: one intentionally simple and one more
realistic. The rules themselves are not claimed as contributions; they
are concrete examples of model use and an introduction to the rule
properties discussed in Section 5.

\begin{itemize}
\tightlist
\item
  \textbf{Example Rule A (simple, reactive):} intent-vs-observation
  correlation after telemetry arrives.
\item
  \textbf{Example Rule B (more realistic, proactive+reactive):}
  pre-action blocking for MCP-gated operations plus Rule A for residual
  detection.
\end{itemize}

\textbf{Example Rule A: Reactive Intent-vs-Behavior Correlation}

\begin{Shaded}
\begin{Highlighting}[]
\NormalTok{process\_allowed = (evidence\_process in allowed\_processes)}
\NormalTok{domain\_allowed  = (target\_ip in allowed\_domains)}
\NormalTok{explained       = process\_allowed AND domain\_allowed}

\NormalTok{guardrail\_signal = (safety\_floor\_findings \textgreater{} 0)}
\NormalTok{                OR (vulnerability\_findings \textgreater{} 0)}

\NormalTok{intent\_mismatch  = (session\_matches \textgreater{} 0 AND explained == false)}
\NormalTok{                OR (credential\_access AND session\_matches \textgreater{} 0)}

\NormalTok{divergence = guardrail\_signal OR intent\_mismatch}
\end{Highlighting}
\end{Shaded}

\textbf{Example Rule B: Pre-Action Guard Rule}

\begin{Shaded}
\begin{Highlighting}[]
\NormalTok{action = proposed tool call (name, arguments)}
\NormalTok{intent = current declared intent}

\NormalTok{spawns\_from\_tmp     = action.working\_dir starts\_with /tmp/ | /var/tmp/ | /dev/shm/}
\NormalTok{targets\_credentials = action.arguments reference sensitive file paths}
\NormalTok{targets\_blacklisted = action.arguments contain known{-}blacklisted endpoints}

\NormalTok{process\_in\_allowlist = action.tool\_binary in intent.allowed\_processes}
\NormalTok{target\_in\_allowlist  = action.target\_endpoint in intent.allowed\_domains}

\NormalTok{static\_violation = spawns\_from\_tmp OR targets\_credentials OR targets\_blacklisted}
\NormalTok{intent\_violation = NOT (process\_in\_allowlist AND target\_in\_allowlist)}

\NormalTok{block = static\_violation OR intent\_violation}
\end{Highlighting}
\end{Shaded}

In this draft, ``sensitive credential paths'' is a configurable static
set (for example \texttt{\textasciitilde{}/.ssh/*}, kube configs, cloud
credential stores, browser-token stores, and CI token files) defined by
deployment policy rather than inferred at runtime.

Rule B does not replace Rule A: actions bypassing MCP still require
reactive system-plane correlation.

\begin{figure}
\centering
\pandocbounded{\includegraphics[keepaspectratio,alt={Detection Decision Flow (Example Rule A).}]{figures/fig4_decision_flow.pdf}}
\caption{Detection Decision Flow (Example Rule A).}
\end{figure}

In this flow, sessions are checked against blacklists,
vulnerability-detection checks, and intent-conditioned allowlists within
a settling-window loop. Divergence is raised at the first failing check.

Model-level multi-signal scenarios (network + identity + LAN + posture)
are provided in Appendix A to keep the main narrative focused.

\subsection{4. Threat Model, Assumptions, and
Boundaries}\label{threat-model-assumptions-and-boundaries}

We focus on runtime compromise modes that evade static checks:

\begin{itemize}
\tightlist
\item
  indirect prompt injection leading to unsafe tool use
\item
  tool/MCP poisoning and capability misuse
\item
  unexplained outbound communications during ostensibly read-only tasks
\item
  credential-path and config-adjacent anomalies
\item
  cross-session behavioral drift and memory poisoning
\end{itemize}

\subsubsection{4.1 Assumptions and Why They Are
Reasonable}\label{assumptions-and-why-they-are-reasonable}

\textbf{Assumptions}

\begin{itemize}
\tightlist
\item
  The observer is co-located with the runtime host.
\item
  Baseline S1/S2 controls are not skipped.
\item
  High-risk actions can require human confirmation when enabled.
\end{itemize}

\textbf{Rationale}

\begin{itemize}
\tightlist
\item
  Co-location is necessary for process attribution and parent-lineage
  fidelity.
\item
  If S1/S2 are absent, runtime detection is overloaded with preventable
  risk.
\item
  Human confirmation for destructive actions is standard in production
  ops.
\end{itemize}

\textbf{Out of scope (with rationale)}

\begin{itemize}
\tightlist
\item
  Formal properties under arbitrarily compromised reasoning planes
  (requires stronger mechanisms such as formal methods or trusted
  execution).
\item
  Internet-scale adaptive adversarial campaigns (requires different
  evaluation infrastructure than the controlled benchmark used here).
\end{itemize}

\subsubsection{4.2 Reasoning-Plane
Integrity}\label{reasoning-plane-integrity}

The key question raised by reviewers is valid: if an attacker can
control intent declarations, can they simply declare harmful intent and
pass correlation? The answer depends on which failure case occurred:

\begin{itemize}
\tightlist
\item
  \textbf{Case A - Tool misuse within declared benign intent.} The model
  is strong. Divergence appears because observed behavior exceeds
  declared intent.
\item
  \textbf{Case B - Correctly declared harmful intent.} Correlation alone
  is insufficient if policy allows that action. Additional governance
  controls (environment-aware allowlists, posture-delta thresholds,
  change-control invariants) are required.
\item
  \textbf{Case C - Missing blast-radius guardrails.} If destructive
  production actions are broadly authorized, the root cause is
  authorization architecture, not telemetry correlation. The two-plane
  model complements S1/S2; it does not replace IAM scoping,
  segmentation, and approval workflows.
\end{itemize}

To raise attacker cost in Cases A/B, we apply two mitigations.

\textbf{Intent-independent critical signals.} The following triggers
cannot be downgraded by LLM adjudication or intent declarations:

\begin{enumerate}
\def\labelenumi{\arabic{enumi}.}
\tightlist
\item
  blacklisted destination access
\item
  CVE-inspired vulnerability-detection failures (gateway binding, token
  exfiltration, skill supply chain, sandbox exploitation)
\end{enumerate}

Model-correlated signals (unexplained sessions, credential access
concurrent with outbound traffic, anomalous network patterns) are
promoted as divergence evidence and adjudicated by the LLM. The LLM
always makes the final verdict decision, but critical signals act as
non-downgradable guardrails.

\textbf{Intent-declaration anomaly tracking.} Motivation: compromised
agents often ``prepare the ground'' by gradually broadening intent
language before executing harmful actions. We therefore track:

\begin{enumerate}
\def\labelenumi{\arabic{enumi}.}
\tightlist
\item
  allowlist budget expansion across redeclarations
\item
  abrupt intent-category shifts
\item
  redeclaration churn (rate spikes)
\end{enumerate}

This signal is currently specified at the model level; its marginal
contribution is not isolated in the Section 7 confusion-matrix metrics.

These mitigations do not remove the trust boundary entirely, but they
force an attacker to simultaneously evade both declaration-level and
system-level checks.

\subsubsection{4.3 Observer Integrity and Evasion
Boundaries}\label{observer-integrity-and-evasion-boundaries}

The model assumes a trustworthy observer process. Under kernel/root
compromise, the observer itself can be tampered with. Potential
hardening paths include watchdog heartbeat, cryptographic attestation,
and stronger isolation.

Primary residual evasion classes are:

\begin{itemize}
\tightlist
\item
  \textbf{Living-off-the-land} (allowed process + allowed destination)
\item
  \textbf{timing attacks} (harm before first informative snapshot)
\item
  \textbf{observer evasion/tampering}
\item
  \textbf{credential-less exfiltration} (transient file access not
  captured in time)
\end{itemize}

The most operationally important combined risk is living-off-the-land
plus timing. This combination is dangerous because behavior can remain
syntactically ``allowed'' while executing before the most informative
telemetry fields are available. Near-term mitigations are shorter
enrichment intervals, sticky process-context retention, stricter
defaults for high-risk tasks, and proactive Rule B gating for
MCP-mediated operations. We detail mitigation priorities and benchmark
blind spots in Section 9.

\subsubsection{4.4 State Isolation and Observer
Independence}\label{state-isolation-and-observer-independence}

A subtle but critical attack surface in agentic frameworks is
\emph{shared mutable state}. Many frameworks persist inter-session
context -- behavioral models, verdicts, operational memory -- in flat
files within the agent's workspace (e.g., \texttt{MEMORY.md},
\texttt{SOUL.md}). These files are readable and writable by any process
on the host. This creates multiple attack vectors:

\begin{enumerate}
\def\labelenumi{\arabic{enumi}.}
\tightlist
\item
  \textbf{Model poisoning.} A compromised process writes false
  behavioral predictions into the shared memory file, causing the
  detection layer to whitelist malicious traffic patterns as
  ``expected.''
\item
  \textbf{Verdict tampering.} A compromised process overwrites the
  stored verdict (e.g., writing \texttt{CLEAN} over a
  \texttt{DIVERGENCE} determination), suppressing alerts before they
  reach human operators.
\item
  \textbf{Self-covering reasoning plane.} A compromised agent can edit
  its own memory to erase evidence of divergent behavior, defeating any
  retrospective audit.
\end{enumerate}

These attacks are particularly dangerous because they require no
privilege escalation -- only filesystem write access to a user-owned
directory.

Our architecture mitigates this by \emph{internalizing detection state}.
The divergence engine runs inside the EDAMAME host runtime: the EDAMAME
Posture daemon in the reference benchmark deployment, and the EDAMAME
Security app on workstation-hosted Cursor paths. Behavioral models,
verdict history, and engine configuration live in process memory and are
accessible only through authenticated MCP endpoints (pre-shared key).
Multiple reasoning-plane producers can push agent-tagged slices via
\texttt{upsert\_behavioral\_model}; EDAMAME merges those slices into one
observer-owned correlation window while retaining contributor
attribution metadata. The reasoning-plane agent can contribute
predictions, but cannot read or modify verdicts, engine state, or
critical-signal rules.

This separation provides three concrete assurances:

\begin{itemize}
\tightlist
\item
  \textbf{Unilateral tamper resistance.} Compromising the reasoning
  plane alone is insufficient to suppress detection. The observer's
  state is unreachable from the agent's process or filesystem.
\item
  \textbf{Independent alerting.} The observer-owned notification
  pipeline (Slack, Telegram, webhook) is shared across agentic and
  divergence/vulnerability outputs, not owned by the agent. Even a fully
  compromised reasoning plane cannot intercept or cancel outbound
  alerts.
\item
  \textbf{Critical-signal invariance.} Intent-independent rules (Section
  4.2) are evaluated inside the observer. No external input can
  downgrade a critical-signal finding.
\end{itemize}

The extrapolator skill retains OpenClaw's memory system for its own
operational checkpoint (a cursor tracking the last-processed session).
This is acceptable: the checkpoint is not security state. Corrupting it
causes the extrapolator to reprocess sessions but cannot influence
EDAMAME's verdicts.

\subsection{5. Properties of the Decision
Rules}\label{properties-of-the-decision-rules}

This section can be read as a short extension of Section 4 (Threat
Model): it states properties of the deterministic decision rules
(Example Rules A-B) that hold by construction. These are not deep
theorems; they are simple, auditable invariants meant to make system
behavior transparent. We explicitly distinguish these from stronger
formal properties (PAC bounds, temporal SMT constraints, cryptographic
proofs) provided by systems such as Pro2Guard, Agent-C, and VET (Section
2.2). Our contribution is practical, not formal.

\subsubsection{5.1 Definitions}\label{definitions}

For each step t:

\begin{itemize}
\tightlist
\item
  I\_t: declared intent (reasoning plane)
\item
  C\_t: context constraints derived from I\_t (e.g., allowlists)
\item
  O\_t: system-plane observation snapshot after settling
\item
  K(O\_t): critical-signal predicate (true if blacklisted traffic or
  vulnerability-detection failure)
\item
  E(I\_t, C\_t, O\_t): ``explained'' predicate (true if O\_t is
  consistent with intent)
\end{itemize}

The two-plane alert predicate is:

\[
D(I_t, C_t, O_t) = K(O_t) \lor (O_t \neq \emptyset \land \neg E(I_t, C_t, O_t))
\]

\subsubsection{5.2 Property 1 (False-Positive
Monotonicity)}\label{property-1-false-positive-monotonicity}

For any anomaly predicate A over observations, let D\_gated = A AND (NOT
E) and D\_critical = K OR D\_gated. Then for benign runs where K does
not fire:

\begin{Shaded}
\begin{Highlighting}[]
\NormalTok{P(D\_critical = 1 | benign, K=0) = P(D\_gated = 1 | benign) \textless{}= P(A = 1 | benign)}
\end{Highlighting}
\end{Shaded}

This follows from set inclusion: \{D\_gated=1\} implies \{A=1\}.
Intent-gating cannot increase false positives relative to anomaly-only
alerting when no critical signal is triggered.

\textbf{Remark (critical-signal false positives).} Critical signals
\emph{can} produce false positives for benign operations that
coincidentally match (e.g., a legitimate connection to a
recently-blacklisted destination).

\subsubsection{5.3 Property 2 (Settling Window and Miss
Rate)}\label{property-2-settling-window-and-miss-rate}

Let L be the time until an action's effects become observable. With
maximum wait W: P(miss) = P(L \textgreater{} W). Progressive re-checks
(sampling at t1 \textless{} t2 \textless{} \ldots{} \textless{} W) do
not change this bound but reduce the expected detection latency when L
\textless{} W.

\subsubsection{5.4 Property 3 (Example Rule A Correctness Under Complete
Telemetry)}\label{property-3-example-rule-a-correctness-under-complete-telemetry}

Assume the evidence bundle is complete:
\texttt{session\_matches\ \textgreater{}\ 0} iff the endpoint was
contacted by \texttt{evidence\_process} after the watermark. Example
Rule A returns \texttt{divergence=true} iff the observed egress is not
authorized by intent OR a critical signal fires. This follows directly
from the algorithm's structure and is not a claim about real-world
completeness. In practice, this assumption is violated by settling
windows and refresh delays (Section 3.6, Section 9.1), so Property 3
should be read as a rule-level sanity property, not an empirical
coverage claim.

\subsubsection{5.5 Property 4 (Pre-Action Blocking is
Conservative)}\label{property-4-pre-action-blocking-is-conservative}

Example Rule B blocks on
\texttt{static\_violation\ OR\ intent\_violation}. Since both predicates
are computed from the proposed action and current intent (no telemetry),
blocking is immediate (zero latency). The cost is over-blocking:
legitimate actions that happen to match a floor pattern (e.g., a build
tool that writes to \texttt{/tmp/}) require human confirmation.

\subsubsection{5.6 Scope of These
Properties}\label{scope-of-these-properties}

These properties characterize the decision rules under idealized
assumptions (complete telemetry, correct allowlists, honest intent for
Property 1). They do not provide:

\begin{itemize}
\tightlist
\item
  Safety properties under partial telemetry (missed sessions, eBPF gaps)
\item
  Bounds on evasion probability (see Section 4.3)
\item
  Temporal safety properties (no ordering constraints on actions)
\item
  Probabilistic properties on detection rates (the Wilson CI in Section
  7.4 is empirical, not analytical)
\end{itemize}

Systems that require these stronger properties should consider Pro2Guard
(PAC bounds), Agent-C (temporal SMT), or VET (cryptographic
attestation).

Part I takeaway: the model provides a portable and auditable way to
correlate declared intent with host-observed behavior, but by itself it
does not provide zero-latency interception, complete living-off-the-land
coverage, or formal safety proofs. Part II therefore focuses on
implementation feasibility and measured operational behavior.

\subsection{Part II --- Implementation and
Evaluation}\label{part-ii-implementation-and-evaluation}

Part II demonstrates feasibility and relevance: Section 6 describes our
implementation, Section 7 reports evaluation results, Section 8
documents reproducibility assets, Section 9 states limitations and
near-term work, and Section 10 discusses ethics and safe use.

\subsection{6. Implementation (OpenClaw + EDAMAME
Posture)}\label{implementation-openclaw-edamame-posture}

Our implementation uses a shared-timer architecture with one OpenClaw
reasoning skill, EDAMAME-internal detection/remediation loops, and an
internal EDAMAME divergence engine:

\textbf{Two-Plane Detection (cooperating timer jobs):}

\begin{itemize}
\tightlist
\item
  \texttt{skill/edamame-extrapolator/} -- Reads raw session history via
  \texttt{sessions\_list(activeMinutes=15)} and
  \texttt{sessions\_history}, distills a compact behavioral model, and
  publishes it via \texttt{upsert\_behavioral\_model}. Runs every 2-5
  minutes. Uses per-session \texttt{message\_count} checkpoints to
  implement a sliding window and avoid re-analysis.
\item
  \textbf{Internal divergence engine} (inside EDAMAME Posture) --
  Implements a hybrid Rust-deterministic + LLM-adjudication
  architecture. Runs every \textasciitilde5 minutes by default
  (configurable). Each tick proceeds in three stages:

  \begin{enumerate}
  \def\labelenumi{\arabic{enumi}.}
  \tightlist
  \item
    \textbf{Deterministic evidence stage} (Rust,
    \texttt{divergence\_engine.rs}): Stateless pure functions correlate
    behavioral model predictions against live telemetry (sessions,
    anomalous, blacklisted, exceptions) and run CVE-inspired
    vulnerability-detection checks (gateway binding via
    \texttt{ourselves.open\_ports}, token exfiltration, skill supply
    chain, sandbox exploitation). Produces a baseline
    \texttt{DivergenceVerdict} with typed evidence (CRITICAL/HIGH/MEDIUM
    severity).
  \item
    \textbf{LLM adjudication stage}: The deterministic evidence bundle
    is submitted to the configured LLM provider for final verdict
    decision and operator-facing alert synthesis. The LLM receives the
    full evidence context and returns a structured
    \texttt{LlmDivergenceDecision} (verdict, confidence, reasoning) plus
    an optional \texttt{LlmHumanAlert} (title, body, severity).
  \end{enumerate}
\end{itemize}

\begin{enumerate}
\def\labelenumi{\arabic{enumi}.}
\setcounter{enumi}{2}
\tightlist
\item
  \textbf{Guardrail reconciliation}
  (\texttt{reconcile\_with\_guardrails}): Critical signals --
  vulnerability-detection findings and blacklisted-traffic evidence --
  are non-downgradable. If the LLM attempts to override a CRITICAL
  deterministic finding, the guardrail forces DIVERGENCE with provenance
  \texttt{GUARDRAIL\_FORCED\_DETERMINISTIC}. Otherwise, the LLM verdict
  is accepted with provenance \texttt{LLM\_CONFIRMED} or
  \texttt{LLM\_OVERRULED\_SOFT\_SIGNALS}. Each verdict records full
  provenance: \texttt{deterministic\_verdict},
  \texttt{decision\_source}, optional \texttt{llm\_decision}, and
  optional \texttt{human\_alert}. If the behavioral model is stale
  (\textgreater20 min), produces a STALE verdict. Verdicts are
  accessible via \texttt{get\_divergence\_verdict} MCP tool.
\end{enumerate}

\textbf{Posture Management (EDAMAME-internal loop):}

\begin{itemize}
\tightlist
\item
  EDAMAME posture/remediation loop (internal) -- Reviews advisor todos
  and executes safe fixes through the native agentic pipeline. In
  OpenClaw, \texttt{skill/edamame-posture/} acts as a thin MCP facade
  for invoking these capabilities; scheduled remediation ownership
  remains inside EDAMAME.
\end{itemize}

\textbf{Supporting components:}

\begin{itemize}
\tightlist
\item
  OpenClaw plugin bridging to MCP: \texttt{extensions/edamame-mcp/}
\item
  Live benchmark harness: \texttt{tests/benchmark/run\_live\_suite.py}
\item
  Metric summarizer: \texttt{tests/benchmark/summarize\_results.sh}
\item
  Canonical claim binding: \texttt{docs/CLAIM\_ARTIFACT\_INDEX.md}
\end{itemize}

Behavioral-model and verdict state are internal to the EDAMAME
divergence engine. OpenClaw memory remains optional operational
checkpointing for the extrapolator cursor and is not authoritative
detection state.

\subsubsection{6.1 MCP Integration Path}\label{mcp-integration-path}

EDAMAME Posture exposes a local, PSK-authenticated MCP endpoint via
Streamable HTTP (\texttt{http://127.0.0.1:3000/mcp}). OpenClaw calls
EDAMAME MCP tools through a local OpenClaw plugin
(\texttt{extensions/edamame-mcp}) that implements a minimal HTTP MCP
client and exposes EDAMAME tools as native OpenClaw tools:

OpenClaw (native tool call) -\textgreater{}\\
\texttt{edamame-mcp} plugin -\textgreater{}\\
EDAMAME MCP (HTTP) -\textgreater{}\\
tool result

\paragraph{MCP Integration Notes}\label{mcp-integration-notes}

Section 3.4 explains why MCP is model-relevant. At implementation level,
the key point is that OpenClaw can consume EDAMAME as native tools
(typed JSON over local PSK-authenticated HTTP), so correlation logic can
run without shell parsing or privileged subprocess glue.

\subsubsection{6.2 Trace-Backed Evidence
Bundles}\label{trace-backed-evidence-bundles}

The live benchmark suite records:

\begin{itemize}
\tightlist
\item
  baseline and post-action telemetry snapshots
\item
  inject logs with VM-side timestamps
\item
  MCP call outputs (todos, action history)
\item
  a run manifest with scenario set hash and git SHA
\end{itemize}

These artifacts enable independent inspection of every verdict.

\subsubsection{6.3 Claim Gating and Validation Maturity
Tiers}\label{claim-gating-and-validation-maturity-tiers}

To avoid confusion between runtime \textbf{evidence bundles} (system
telemetry used by the divergence engine) and paper-level
\textbf{validation status} (what this manuscript has empirically
established), we use the following maturity tiers:

\begin{itemize}
\tightlist
\item
  T0: implemented
\item
  T1: demonstrated in integration tests
\item
  T2: measured in trace-backed benchmark artifacts
\item
  T3: independently replicated
\end{itemize}

Automation blocks unsupported strong-claim language unless minimum tier
floors are met. Published metric claims are anchored to a single
manifest/summary lineage in \texttt{docs/CLAIM\_ARTIFACT\_INDEX.md}, so
run IDs, hashes, and benchmark modes remain auditable.

\subsubsection{6.4 Signal Coverage vs Validation Tier (What Is
Empirically
Supported)}\label{signal-coverage-vs-validation-tier-what-is-empirically-supported}

The system targets multi-signal attacks, but empirical support is
non-uniform. Table 1 summarizes current status:

{\def\LTcaptype{none} % do not increment counter
\begin{longtable}[]{@{}
  >{\raggedright\arraybackslash}p{(\linewidth - 4\tabcolsep) * \real{0.4091}}
  >{\raggedright\arraybackslash}p{(\linewidth - 4\tabcolsep) * \real{0.2727}}
  >{\raggedright\arraybackslash}p{(\linewidth - 4\tabcolsep) * \real{0.3182}}@{}}
\toprule\noalign{}
\begin{minipage}[b]{\linewidth}\raggedright
Channel
\end{minipage} & \begin{minipage}[b]{\linewidth}\raggedright
Tier
\end{minipage} & \begin{minipage}[b]{\linewidth}\raggedright
Notes
\end{minipage} \\
\midrule\noalign{}
\endhead
\bottomrule\noalign{}
\endlastfoot
Network egress (sessions, anomalous, blacklisted) & T2 & Trace-backed
live suite (paper metrics) \\
Process lineage (parent\_pid, parent\_name, parent\_path) & T1 & Layer-7
enrichment integration-tested via MCP \\
Open file access (per-process file descriptors) & T1 & Credential-theft
detection via open\_files \\
Exceptions / policy violations & T1 & Sanity-checked, no dedicated T2
protocol \\
MCP wiring (advisor\_get\_todos, action\_history) & T1 & OpenClaw
-\textgreater{} plugin -\textgreater{} MCP path verified \\
Posture scoring (get\_score, advisor\_get\_todos) & T1 & Triage/demo
flows; not yet T2-measured \\
Vulnerability detection (config/env safety scans) & T1 & Unsafe
configuration plus host-open-port gateway exposure checks \\
Cross-session drift (memory poisoning) & T0/T1 & Partially exercised;
not in T2 protocol \\
\end{longtable}
}

\subsubsection{6.5 Operational Context: Posture Scoring and
Compliance}\label{operational-context-posture-scoring-and-compliance}

Beyond individual network signals, EDAMAME Posture continually computes
a composite security score (0-100) with subcategory breakdowns (network,
system integrity, system services, applications, credentials) and maps
the endpoint state against compliance frameworks (CIS Benchmark Level
1/2, ISO 27001/2, SOC 2). This provides two advantages for the two-plane
model:

\begin{enumerate}
\def\labelenumi{\arabic{enumi}.}
\item
  \textbf{Baseline context.} Endpoint hardening level (e.g., 90/100 vs
  45/100) changes risk posture and response aggressiveness.
\item
  \textbf{Score-delta detection.} A posture drop during a task (e.g., 65
  to 50) is an independent system-plane signal, even when individual
  sessions look benign.
\end{enumerate}

These capabilities are exercised in the extrapolator/internal divergence
engine flow and demo path (T1). Their incremental impact on T2 live
detection accuracy remains future work. Even at current maturity, this
context already improves operator triage and supports safer action
review in T1-validated workflows.

\subsection{7. Evaluation: The BadAgentUse
Benchmark}\label{evaluation-the-badagentuse-benchmark}

We introduce \textbf{BadAgentUse}, an open benchmark for evaluating
agentic-system monitors against real endpoint telemetry. Unlike prior
benchmarks that simulate network traces or rely on sandboxed
environments, BadAgentUse injects live traffic into a production-grade
Linux VM and evaluates detection against actual EDAMAME Posture session
captures.

\subsubsection{7.1 Setup}\label{setup}

\begin{itemize}
\tightlist
\item
  \textbf{Environment:} Lima VM (\texttt{openclaw-security}), Ubuntu
  24.04, 4 vCPU / 4 GiB
\item
  \textbf{Observer:} EDAMAME Posture daemon (eBPF-backed on aarch64
  Linux) with MCP Streamable HTTP endpoint on port 3000
\item
  \textbf{Reasoning-plane path:} In production, OpenClaw publishes
  behavioral-model slices through the \texttt{edamame-extrapolator}
  skill, and Cursor publishes through the same EDAMAME contributor
  contract. For the quantitative BadAgentUse runs reported here,
  however, the harness injects scenario-specific behavioral-model slices
  directly via \texttt{upsert\_behavioral\_model} after deriving them
  from the scenario contract and live baseline telemetry. This isolates
  the downstream divergence engine from cron cadence and LLM
  nondeterminism. End-to-end extrapolator behavior is validated
  separately by integration and E2E tests rather than by the main
  aggregate metric loop.
\item
  \textbf{Protocol:} 50 versioned JSON scenario files in
  \texttt{tests/benchmark/live-scenarios/}
\item
  \textbf{Iterations:} configurable via \texttt{-\/-iterations}; the
  current canonical artifact lineage is a partial live sweep rather than
  a uniform full grid
\item
  \textbf{Seeds:} 12 distinct seeds in the current canonical live
  evidence set, with deterministic target-IP rotation
\item
  \textbf{Trial unit:} one scenario evaluated per seed-iteration pair.
  The confusion matrix (Section 7.4) reports n=273 valid runs across 12
  seeds, with 172 skipped runs preserved in the readiness scorecard.
  Each iteration varies seeds and settling behavior; all iteration
  traces remain in artifacts.
\end{itemize}

\paragraph{Protocol Parameters}\label{protocol-parameters}

{\def\LTcaptype{none} % do not increment counter
\begin{longtable}[]{@{}
  >{\raggedright\arraybackslash}p{(\linewidth - 4\tabcolsep) * \real{0.3333}}
  >{\raggedright\arraybackslash}p{(\linewidth - 4\tabcolsep) * \real{0.3333}}
  >{\raggedright\arraybackslash}p{(\linewidth - 4\tabcolsep) * \real{0.3333}}@{}}
\toprule\noalign{}
\begin{minipage}[b]{\linewidth}\raggedright
Parameter
\end{minipage} & \begin{minipage}[b]{\linewidth}\raggedright
Value
\end{minipage} & \begin{minipage}[b]{\linewidth}\raggedright
Notes
\end{minipage} \\
\midrule\noalign{}
\endhead
\bottomrule\noalign{}
\endlastfoot
Settling schedule & 15 s → 30 s → 60 s & Exponential back-off; max wait
W = 60 s \\
Stability criterion & 2 consecutive stable snapshots & Snapshot =
sessions + Layer-7 + posture unchanged \\
Layer-7 full refresh & 300 s (5 min) & Process metadata, parent lineage,
open files \\
Sensitive-file scan & 30 s & Credential paths only; sticky retention \\
Anomaly-score threshold & 0 (any non-zero triggers) & Conservative;
tunable per deployment \\
Posture-delta threshold & \(\geq\) 1 point drop & Monitoring signal (not
hard-gate) \\
Blacklist sources & FireHOL Level 1 + custom & Pinned at VM provision
time \\
Allowlist semantics & Exact domain OR IP; no wildcards & AND logic:
process AND destination must match \\
MCP endpoint & \texttt{127.0.0.1:3000}, PSK-authenticated &
Localhost-only by default \\
LLM (attack detection) & \texttt{openai/gpt-5.1} & Deterministic tiers
do not use LLM; compatible Azure-hosted OpenAI Responses models are
supported as alternate providers \\
\end{longtable}
}

\subsubsection{7.2 Scenario Design (Versioned
Categories)}\label{scenario-design-versioned-categories}

In BadAgentUse, a \textbf{scenario} is a versioned JSON test contract
that specifies: (1) declared intent context, (2) injected runtime action
or traffic pattern, (3) expected enforcement label (\texttt{DIVERGENCE},
\texttt{CLEAN}, or \texttt{DESIGNED\_MISS}), and (4) the telemetry
checks needed for scoring.

BadAgentUse covers eight categories spanning five signal families:

{\def\LTcaptype{none} % do not increment counter
\begin{longtable}[]{@{}
  >{\raggedright\arraybackslash}p{(\linewidth - 6\tabcolsep) * \real{0.2000}}
  >{\raggedright\arraybackslash}p{(\linewidth - 6\tabcolsep) * \real{0.1000}}
  >{\raggedright\arraybackslash}p{(\linewidth - 6\tabcolsep) * \real{0.2667}}
  >{\raggedright\arraybackslash}p{(\linewidth - 6\tabcolsep) * \real{0.4333}}@{}}
\toprule\noalign{}
\begin{minipage}[b]{\linewidth}\raggedright
Cat.
\end{minipage} & \begin{minipage}[b]{\linewidth}\raggedright
N
\end{minipage} & \begin{minipage}[b]{\linewidth}\raggedright
Signal
\end{minipage} & \begin{minipage}[b]{\linewidth}\raggedright
Description
\end{minipage} \\
\midrule\noalign{}
\endhead
\bottomrule\noalign{}
\endlastfoot
A. Network-Intent & 17 & Network & Undeclared egress, partial-allowlist
AND, lateral movement \\
B. Process Lineage & 6 & Layer-7 & /tmp/ scripts, interpreter wrap,
nested chains \\
C. Credential Theft & 7 & Layer-7 files & SSH, kube, gnupg, browser
cookies, combined exfil \\
D. CVE Vulnerability Detection & 6 & Layer-7+allow & gateway\_binding /
token\_exfil / credential overrides \\
E. Posture \& Policy & -- & Posture & Score, policy, threats, HIBP
(integration checks)† \\
F. LAN Awareness & -- & LAN & Device discovery, scan, anomaly APIs
(integration checks)† \\
G. Adversarial & 8 & Mixed & LotL, timing, burst, path substitution \\
H. Noisy Benign & 6 & Network & Clean traffic, DNS, NTP, auto-update \\
\end{longtable}
}

† Categories E and F are validated via dedicated integration scripts
(Section 7.7), not as JSON benchmark scenarios. The 50 benchmark
scenarios span categories A--D, G, H.

\textbf{Enforcement labels.} Each scenario carries an
\texttt{expected\_enforcement} field: \texttt{DIVERGENCE} (39
scenarios), \texttt{CLEAN} (8 scenarios), or \texttt{DESIGNED\_MISS} (3
scenarios). The \texttt{DESIGNED\_MISS} label explicitly marks
living-off-the-land attacks that Example Rule A cannot detect by
construction (Section 4.3).

\subsubsection{7.3 Metrics}\label{metrics}

We adopt the metric framework from AgentSentinel {[}Hu et al., 2025{]}:

\begin{itemize}
\tightlist
\item
  \textbf{DSR} (Detection Success Rate) = TP / (TP + FN) -- equivalent
  to recall for attack scenarios
\item
  \textbf{FPR} (False Positive Rate) = FP / (FP + TN) -- rate of benign
  traffic incorrectly flagged
\item
  \textbf{FNR} (False Negative Rate) = FN / (FN + TP) -- attacks that
  evaded detection
\item
  \textbf{ASR} (Attacker Success Rate) = FNR -- attacks that succeeded
  despite monitoring
\item
  \textbf{FT} (False Trigger) = FPR -- identical to FPR for symmetry
\item
  \textbf{TTD} (Time-to-Detect) = median and p95 detection latency (ms),
  computed from injection timestamp to first evidence capture
\item
  \textbf{Wilson CI95} on DSR and FPR for statistical significance
\end{itemize}

All metrics are computed from trace-backed NDJSON results by
\texttt{summarize\_results.sh}, which also produces per-category and
per-seed breakdowns for stability analysis.

\subsubsection{7.4 Statistical Confidence Reporting (Wilson
Interval)}\label{statistical-confidence-reporting-wilson-interval}

To bound measured proportions (precision/recall), we use the Wilson
score interval (Wilson, 1927). For \texttt{k} successes out of
\texttt{n} trials and \texttt{z=1.96} (95\% confidence):

\begin{Shaded}
\begin{Highlighting}[]
\NormalTok{p\_hat = k/n}
\NormalTok{den   = 1 + z\^{}2/n}
\NormalTok{center= (p\_hat + z\^{}2/(2n)) / den}
\NormalTok{half  = z * sqrt(p\_hat*(1{-}p\_hat)/n + z\^{}2/(4n\^{}2)) / den}
\NormalTok{CI95  = [center {-} half, center + half]}
\end{Highlighting}
\end{Shaded}

In our current canonical live evidence set:

{\def\LTcaptype{none} % do not increment counter
\begin{longtable}[]{@{}lll@{}}
\toprule\noalign{}
& Predicted DIVERGENCE & Predicted CLEAN \\
\midrule\noalign{}
\endhead
\bottomrule\noalign{}
\endlastfoot
\textbf{Ground-truth attack} & TP = 225 & FN = 3 \\
\textbf{Ground-truth benign} & FP = 36 & TN = 9 \\
\end{longtable}
}

Precision = 225/261 = 0.862 and recall = 225/228 = 0.987, yielding CI95
{[}0.815, 0.899{]} and {[}0.962, 0.996{]}, respectively (n = 273 valid
runs, 12 seeds).

The 3 FN are highly concentrated in two stress buckets: 2 occur in
adversarial-evasion behavior, with the remaining miss in
intent-proof-floor timing. Likewise, 22 of the 36 FP arise from
noisy-benign traffic, with the remaining 14 in the broader
network-intent-correlation bucket. The aggregate result therefore
reflects localized boundary conditions rather than uniform degradation
across categories, and points directly to two remaining needs: stronger
benign-traffic baselining and better coverage of living-off-the-land
behavior.

\subsubsection{7.5 Results}\label{results}

\paragraph{Current Canonical Live Evidence
Set}\label{current-canonical-live-evidence-set}

\emph{Results below come from the canonical evidence set in
\texttt{artifacts/live-paper-summary.json} (50 planned scenarios; 273
valid runs across 12 seeds, with 172 skipped runs retained in the
readiness scorecard). These aggregate metrics are from the
OpenClaw-derived live benchmark population under the explicit-slice
producer mode (\texttt{upsert\_behavioral\_model}) into EDAMAME; Cursor
uses the raw-session producer mode
(\texttt{upsert\_behavioral\_model\_from\_raw\_sessions}) but is not a
separate benchmark cohort in this draft.}

\textbf{Claim-binding metadata (canonical evidence lineage):}

\begin{itemize}
\tightlist
\item
  \texttt{run\_id}: \texttt{run-live-20260304T234441Z-live}
\item
  \texttt{git\_sha}: \texttt{473465c}
\item
  \texttt{scenario\_set\_version}:
  \texttt{22cb4a6cbddf88154435207c279143f7245de94dbc47c5803ca349eb10f3de1d}
\item
  \texttt{mode}: \texttt{live}
\item
  \texttt{benchmark\_mode}: \texttt{live}
\end{itemize}

{\def\LTcaptype{none} % do not increment counter
\begin{longtable}[]{@{}ll@{}}
\toprule\noalign{}
Metric & Value \\
\midrule\noalign{}
\endhead
\bottomrule\noalign{}
\endlastfoot
Total runs & 273 \\
Precision & 0.8620689655172413 \\
Recall & 0.9868421052631579 \\
Precision CI95 & {[}0.8149560047336936, 0.898678107127887{]} \\
Recall CI95 & {[}0.9620350252457207, 0.9955153022192381{]} \\
Median latency (ms) & 45810 \\
p95 latency (ms) & 53582 \\
Rollback reliability & 1 \\
Stability (precision stddev) & 0.10530453269195077 \\
Stability (recall stddev) & 0.07229988054812211 \\
\end{longtable}
}

\textbf{Readiness note.} The current committed lineage remains below the
repo's publication gate:
\texttt{artifacts/arxiv-readiness-scorecard.json} is still
\texttt{NO\_GO} because the precision CI95 lower bound is 0.815 and the
skipped-run ratio is 38.7\%. We keep these values in the draft because
they are the currently reproducible canonical artifacts on this branch.

\paragraph{Per-Category Detection}\label{per-category-detection}

{\def\LTcaptype{none} % do not increment counter
\begin{longtable}[]{@{}
  >{\raggedright\arraybackslash}p{(\linewidth - 8\tabcolsep) * \real{0.2941}}
  >{\raggedright\arraybackslash}p{(\linewidth - 8\tabcolsep) * \real{0.2059}}
  >{\raggedright\arraybackslash}p{(\linewidth - 8\tabcolsep) * \real{0.1471}}
  >{\raggedright\arraybackslash}p{(\linewidth - 8\tabcolsep) * \real{0.1471}}
  >{\raggedright\arraybackslash}p{(\linewidth - 8\tabcolsep) * \real{0.2059}}@{}}
\toprule\noalign{}
\begin{minipage}[b]{\linewidth}\raggedright
Category
\end{minipage} & \begin{minipage}[b]{\linewidth}\raggedright
Run N
\end{minipage} & \begin{minipage}[b]{\linewidth}\raggedright
DSR
\end{minipage} & \begin{minipage}[b]{\linewidth}\raggedright
FPR
\end{minipage} & \begin{minipage}[b]{\linewidth}\raggedright
Notes
\end{minipage} \\
\midrule\noalign{}
\endhead
\bottomrule\noalign{}
\endlastfoot
Network intent correlation & 68 & 1.000 & 0.700 & undeclared egress is
covered well, but benign allowance ambiguity still produces many false
triggers \\
Process lineage verification & 45 & 1.000 & -- & spawned\_from\_tmp,
wrapper lineage, nested chains \\
Credential theft detection & 36 & 1.000 & -- & SSH, kube, gnupg,
browser-cookie access patterns \\
Credential-network combined & 7 & 1.000 & -- & combined credential +
egress behavior \\
Intent-proof floor & 37 & 0.973 & -- & one remaining miss in
floor-policy timing \\
Lateral movement & 1 & 1.000 & -- & nc-based lateral-probe preparation
behavior \\
Adversarial evasion & 54 & 0.963 & -- & most cases detect; two
LotL/timing misses remain \\
Noisy benign & 25 & -- & 0.880 & dominant benign false-positive source
in the current live set \\
\end{longtable}
}

Categories E (Posture \& Policy) and F (LAN Awareness) are validated via
integration scripts with 100\% pass rate (Section 7.7).

\paragraph{Analysis}\label{analysis}

\textbf{Observed false negatives (3).} Most misses occur in the
adversarial-evasion bucket (2 of 3), with one remaining miss in
intent-proof-floor timing. That concentration is more informative than
the raw count: the system remains strong on process-lineage,
credential-access, combined credential-egress, lateral-movement, and
core network-intent categories, while the residual miss surface is
localized to living-off-the-land executions that preserve expected
binaries/destinations and to cases where Layer-7 context arrives after
the initial decision window.

\textbf{Layer-7 timing sensitivity.} The residual intent-proof-floor
miss still looks like an observability-boundary issue rather than a
rule-construction failure. When the 5-minute Layer-7 full-refresh cycle
has not yet captured process-lineage context at snapshot time, the first
decision can be too early; once enrichment lands, the same underlying
logic succeeds.

\textbf{Operational false positives (benign ambiguity).} False positives
are still localized rather than uniform: 22 of the 36 FP arise in the
noisy-benign bucket (\texttt{FPR=0.880}, 22/25 runs), with the remaining
14 in the broader network-intent-correlation category. This suggests
that the current burden is split between background traffic and
benign-but-under-specified allowlist contexts rather than generalized
instability across all categories. In deployment, that still matters:
persistent benign alerts erode operator trust quickly, so
environment-specific allowlisting or baseline learning is needed before
policy hardening.

High-level takeaway: the current 12-seed live evidence set reinforces
the system's value as a runtime correlation layer for lineage,
credential access, lateral movement, and undeclared egress, while
sharpening two caveats. First, benign background traffic still needs
better baseline treatment to avoid operator fatigue. Second,
adversarial/LotL behavior that stays within expected tools or appears
before Layer-7 enrichment remains difficult. The system is therefore
best understood as a practical early-warning layer, not as an
interception-grade prevention mechanism.

\subsubsection{7.6 Ablation Studies}\label{ablation-studies}

To interpret the overall Section 7.5 result, we run ablations that
isolate which components produce most of the detection lift and where
misses originate. The benchmark runner supports three ablation modes via
\texttt{-\/-ablation}:

{\def\LTcaptype{none} % do not increment counter
\begin{longtable}[]{@{}
  >{\raggedright\arraybackslash}p{(\linewidth - 4\tabcolsep) * \real{0.2703}}
  >{\raggedright\arraybackslash}p{(\linewidth - 4\tabcolsep) * \real{0.2703}}
  >{\raggedright\arraybackslash}p{(\linewidth - 4\tabcolsep) * \real{0.4595}}@{}}
\toprule\noalign{}
\begin{minipage}[b]{\linewidth}\raggedright
Ablation
\end{minipage} & \begin{minipage}[b]{\linewidth}\raggedright
Disabled
\end{minipage} & \begin{minipage}[b]{\linewidth}\raggedright
Expected Impact
\end{minipage} \\
\midrule\noalign{}
\endhead
\bottomrule\noalign{}
\endlastfoot
no-introspection (ablation flag name) & CVE vulnerability detection &
Cat. D degrades to CLEAN \\
no-l7 & Layer-7 enrichment & Cat. B, C lose detection \\
no-settling & Settling window & Transient misses; timing evasion \\
\end{longtable}
}

These ablations define component contributions. Running each mode writes
comparative summaries to \texttt{artifacts/live-paper-summary.json} for
the benchmark-summary visualization used in companion materials. In this
draft, full \texttt{no-floor} / \texttt{no-l7} / \texttt{no-settling}
comparative tables are not yet included in the canonical
confusion-matrix snapshot.

\textbf{eBPF vs netstat2 fallback.}
\texttt{tests/test\_l7\_ebpf\_ablation.sh} runs the same Layer-7 checks
twice: with eBPF kprobes (default on kernels \(\geq\) 5.8), then with
tracefs unmounted to force netstat2 \texttt{/proc/net/*} fallback.

{\def\LTcaptype{none} % do not increment counter
\begin{longtable}[]{@{}lllll@{}}
\toprule\noalign{}
Mode & Checks & Pass & Layer-7 Latency & Sessions \\
\midrule\noalign{}
\endhead
\bottomrule\noalign{}
\endlastfoot
eBPF (kprobe) & 9 & 9/9 & 15 s & 50 \\
netstat2 (fallback) & 9 & 9/9 & 45 s & 31 \\
\end{longtable}
}

Both modes pass 9/9 and produce the same enrichment fields
(\texttt{pid}, \texttt{process\_name}, parent lineage including
\texttt{parent\_script\_path}, \texttt{spawned\_from\_tmp},
\texttt{open\_files}, cwd, memory, cpu\_usage). Latency differs (15 s
with eBPF vs 45 s fallback) because eBPF only accelerates socket-to-PID
mapping. After PID resolution, both paths share the same \texttt{/proc}
enrichment logic. Fallback can miss very short-lived sockets between
polls, but for connections that survive at least one poll cycle,
detection capability is preserved.

\subsubsection{7.7 Multi-Signal Observer
Validation}\label{multi-signal-observer-validation}

After isolating component-level effects, we verify observer breadth: the
system plane is not only a network monitor, but a multi-signal telemetry
surface. Categories E (Posture) and F (LAN) validate that the
system-plane observer extends beyond network sessions. The test suite
verifies MCP tool accessibility for five signal families:

\begin{enumerate}
\def\labelenumi{\arabic{enumi}.}
\tightlist
\item
  \textbf{Network sessions:} \texttt{get\_sessions},
  \texttt{get\_exceptions} (Categories A-D, G-H)
\item
  \textbf{Anomaly/blacklist:} \texttt{get\_anomalous\_sessions},
  \texttt{get\_blacklisted\_sessions}
\item
  \textbf{Posture:} \texttt{get\_score}, \texttt{check\_policy},
  \texttt{advisor\_get\_todos} (Category E)
\item
  \textbf{LAN:} \texttt{get\_lan\_devices},
  \texttt{get\_lan\_host\_device}, \texttt{set\_lan\_auto\_scan}
  (Category F)
\item
  \textbf{Identity:} \texttt{get\_pwned\_emails}, \texttt{get\_breaches}
  (Category E)
\end{enumerate}

For benchmark/tool-access validation, the suite calls
\texttt{get\_sessions}; the internal divergence and vulnerability loops
themselves consume \texttt{get\_current\_sessions} within the EDAMAME
host runtime.

All five families are exercised deterministically (no LLM) with 100\%
pass rates, confirming that the two-plane correlation loop can
incorporate signals from any of these dimensions without infrastructure
changes.

\subsubsection{7.8 Empirical Validation: Four-Tier Test
Architecture}\label{empirical-validation-four-tier-test-architecture}

Beyond the benchmark suite, we maintain a four-tier test architecture:

\textbf{Tier 0: Smoke and component validation} (\texttt{smoke.sh},
\texttt{test\_skill\_components.sh},
\texttt{test\_extrapolator\_validation.sh}, \texttt{test\_posture.sh}).
Deterministic checks for daemon health, MCP tool visibility, session
APIs, behavioral model engine handoff, posture skill runtime. 79+ checks
across 4 test groups, all passing.

\textbf{Tier 1: Layer-7 telemetry validation}
(\texttt{test\_l7\_telemetry.sh}, \texttt{test\_l7\_ebpf\_ablation.sh}).
Deterministic checks for parent lineage, open-files detection, field
correctness, and coverage (20 checks plus 2x9 eBPF/fallback checks).

\textbf{Tier 2: Attack and policy/LAN detection}
(\texttt{test\_attack\_detection.sh}, \texttt{test\_posture\_policy.sh},
\texttt{test\_lan\_awareness.sh}). Engine-driven scenarios use the
hybrid Rust-deterministic + LLM-adjudication pipeline for attack
detection, alongside deterministic policy and LAN checks to validate
posture/todos and LAN/anomaly/blacklist/HIBP APIs.

\textbf{Tier 3: Two-plane E2E and vulnerability detection}
(\texttt{test\_two\_plane\_e2e.sh},
\texttt{test\_vulnerability\_detection.sh}). Full worker + extrapolator
+ inject + detector pipeline validation plus CVE-inspired
vulnerability-detection checks across token exfiltration, exposed
gateways, weaponized skills, sandbox escapes, credential sprawl, and
memory poisoning.

\begingroup\footnotesize
\begin{center}
\setlength{\tabcolsep}{4pt}
\renewcommand{\arraystretch}{1.08}
\begin{tabular}{@{}>{\raggedright\arraybackslash}p{0.05\linewidth}>{\raggedright\arraybackslash}p{0.24\linewidth}>{\raggedright\arraybackslash}p{0.07\linewidth}>{\raggedright\arraybackslash}p{0.08\linewidth}>{\raggedright\arraybackslash}p{0.17\linewidth}>{\raggedright\arraybackslash}p{0.27\linewidth}@{}}
\toprule
Tier & Script & Checks & LLM? & Pass Rate & Notes \\
\midrule
0 & \texttt{test\_poc.sh} & 29 & Mixed & 97\% (28/28 + 1 skip) & Infrastructure, telemetry, MCP tools, agent \\
0 & \texttt{smoke.sh} & 14 & No & 100\% & Daemon health, eBPF kprobes \\
0 & \shortstack[l]{\texttt{test\_skill\_}\\\texttt{components.sh}} & 55 & No & 100\% & MCP tools, engine handoff \\
0 & \shortstack[l]{\texttt{test\_extrapolator\_}\\\texttt{validation.sh}} & -- & Yes & 100\% & Behavioral model pipeline \\
0 & \texttt{test\_posture.sh} & 24 & Yes & 100\% & Skill runtime invocation \\
1 & \shortstack[l]{\texttt{test\_l7\_}\\\texttt{telemetry.sh}} & 20 & No & 100\%* & PID, process name, cmd attribution \\
1 & \shortstack[l]{\texttt{test\_l7\_ebpf\_}\\\texttt{ablation.sh}} & 2x9 & No & Partial & eBPF vs fallback comparison \\
2 & \shortstack[l]{\texttt{test\_attack\_}\\\texttt{detection.sh}} & 22 & Mixed & Partial** & Engine verdict pipeline under stabilization \\
2 & \shortstack[l]{\texttt{test\_posture\_}\\\texttt{policy.sh}} & 8 & No & 87\% & 7/8 pass \\
2 & \shortstack[l]{\texttt{test\_lan\_}\\\texttt{awareness.sh}} & 6 & No & 83\% & 5/6 pass \\
3 & \shortstack[l]{\texttt{test\_two\_plane\_}\\\texttt{e2e.sh}} & -- & Yes & Partial** & Worker + extrapolator + detector pipeline \\
3 & \shortstack[l]{\texttt{test\_vulnerability\_}\\\texttt{detection.sh}} & 66 & Yes & 86\% & 57/66 pass \\
\bottomrule
\end{tabular}
\end{center}
\endgroup

*L7 Tier 1: Basic L7 attribution (PID, process\_name, full command line)
passes consistently with eBPF kprobe programs active
(\texttt{track\_connect\_v4/v6}, \texttt{trace\_udp\_send},
\texttt{trace\_sendto}). Extended enrichment fields
(parent\_process\_path, open\_files) are populated asynchronously by the
L7 enricher and depend on \texttt{/proc} scan depth; these fields are
not tested in the core validation suite since they are additive context
for the divergence engine rather than structural correctness
requirements.

**Engine verdict: The hybrid Rust/LLM divergence engine returns
structured verdicts only after a behavioral model has been ingested and
the 60-second settling window has elapsed. The \texttt{test\_poc.sh}
suite correctly handles this as a SKIP condition when no agentic tick
has occurred. Tier 0 component tests (55/55) confirm all engine MCP
tools are reachable.

\subsubsection{7.9 Discussion}\label{discussion}

\textbf{Latency is telemetry settling, not model delay.} The observed
40-60s TTD is dominated by EDAMAME session aggregation and Layer-7
refresh cadence. Without settling-window checks (Section 3.6), many runs
would be scored on stale data.

\textbf{Process attribution is load-bearing.} Attack and benign runs can
generate identical traffic to the same endpoint; only process
attribution separates \texttt{curl} from \texttt{python3} at the same
IP. Without Layer-7 enrichment, separation collapses.

\textbf{Claim discipline is mandatory.} We do not claim comprehensive
coverage. BadAgentUse validates five signal families and explicitly
labels three designed blind spots. The suite is extensible via new JSON
scenarios, without code changes.

\textbf{Operational burden of benign noise.} The noisy-benign slice
currently shows \texttt{FPR=0.40} (2/5). This is acceptable for
controlled benchmarking but high for always-on production operations;
reducing benign-noise alerts is a deployment priority alongside
improving attack recall.

\textbf{Comparative baselines are currently absent.} BadAgentUse does
not yet include head-to-head comparisons against AgentSentinel,
AgentArmor, LlamaFirewall, Pro2Guard, or Agent-C. This benchmark
validates internal correctness of the correlation protocol and telemetry
surface, not superiority. Direct comparison requires adapting either the
harness for interception-style monitors or those systems for MCP
telemetry (Section 9.2).

\subsection{8. Reproducibility}\label{reproducibility}

All source code, scenario definitions, test scripts, and evaluation
harnesses are published at:
\textbf{https://github.com/edamametechnologies/openclaw\_security}

Key artifact paths:

\begin{itemize}
\tightlist
\item
  \texttt{tests/benchmark/live-scenarios/}: versioned JSON scenario
  contracts for live runs
\item
  \texttt{tests/benchmark/run\_live\_suite.py}: benchmark harness
\item
  \texttt{tests/benchmark/summarize\_results.sh}: metric computation
  from NDJSON traces
\item
  \texttt{reproduce\_live.sh}: single-command full reproduction
\item
  \texttt{artifacts/}: generated results (summary JSON, NDJSON traces,
  paper bundle)
\item
  \texttt{artifacts/exec-demo-pack-\textless{}timestamp\textgreater{}/}:
  latest 4-scenario executive demo artifact pack (MP4 + verification
  contracts)
\item
  \texttt{docs/CLAIM\_ARTIFACT\_INDEX.md}: canonical claim-to-artifact
  lineage binding
\end{itemize}

Single-command reproduction for paper metrics:

\begin{Shaded}
\begin{Highlighting}[]
\ExtensionTok{./reproduce\_live.sh}
\end{Highlighting}
\end{Shaded}

This generates:

\begin{itemize}
\tightlist
\item
  \texttt{artifacts/live-paper-results.ndjson}
\item
  \texttt{artifacts/live-paper-summary.json}
\item
  \texttt{artifacts/live-paper-manifest.json}
\item
  \texttt{artifacts/paper-bundle-live/*} (generated whitepaper bundle
  with embedded metrics)
\end{itemize}

\subsection{9. Limitations and Future
Work}\label{limitations-and-future-work}

\subsubsection{9.1 Current Limitations}\label{current-limitations}

\begin{enumerate}
\def\labelenumi{\arabic{enumi}.}
\tightlist
\item
  \textbf{Living-off-the-land blind spot.} The primary residual risk is
  LotL behavior (Section 4.3): allowed processes exfiltrating to allowed
  destinations. BadAgentUse explicitly labels 3 of 50 scenarios as
  \texttt{DESIGNED\_MISS} for this reason. Closing this gap requires
  behavioral baselines, anomaly fusion, or payload-aware inspection.
\item
  \textbf{Layer-7 refresh latency.} Full Layer-7 enrichment refresh runs
  every 5 minutes (sensitive-file scan every 30 s on Linux). Correlation
  rules that rely on \texttt{spawned\_from\_tmp} or \texttt{open\_files}
  can miss short-lived processes if snapshots are stale. The
  \texttt{no-settling} ablation (Section 7.6) quantifies this.
\item
  \textbf{Reactive detection delay (operationally significant).} Example
  Rule A is post-observation and therefore delayed (median 45.8 s). In
  high-consequence environments, this window may still allow harmful
  actions before alerting. Example Rule B reduces this for MCP-gated
  actions, but direct shell/library paths can bypass pre-action checks.
\item
  \textbf{Interception trade-off vs AgentSentinel.} AgentSentinel
  reduces timing exposure for intercepted operations by pausing
  execution inline; our approach prioritizes portability and lower
  integration friction. This is a real design trade-off, not a
  superiority claim.
\item
  \textbf{Governance boundary (Cases B/C).} If the reasoning plane
  truthfully declares a harmful action and policy allows it, correlation
  alone is insufficient (Section 4.2). Authorization architecture,
  change control, and approvals remain mandatory.
\item
  \textbf{Reasoning-plane observation (partially addressed).} The
  extrapolator reads session transcripts via platform APIs, providing
  independent observation of tool calls and arguments. However, this is
  transcript-only -- chain-of-thought reasoning, internal model state,
  and unlogged deliberation are not observable. A sufficiently
  sophisticated attacker could craft tool calls that appear benign in
  the transcript while achieving malicious outcomes through side effects
  invisible to the session record.
\item
  \textbf{Observer integrity boundary.} Section 4.3 assumes a
  trustworthy observer. Under kernel-level compromise, telemetry can be
  suppressed or forged.
\item
  \textbf{Single-platform benchmark.} BadAgentUse is currently
  Linux-only (aarch64 Lima VM). Multiplatform support exists in the
  observer but is not yet fully benchmarked in this protocol.
\item
  \textbf{Single-observer, non-federated focus.} The current EDAMAME
  runtime already supports multiple reasoning-plane contributors on one
  observer-hosted engine (for example OpenClaw plus Cursor slices merged
  into one correlation window), but multi-host delegation still requires
  federated observation and cross-host correlation that are not yet
  implemented.
\item
  \textbf{No formal safety proofs.} Section 5 provides transparent rule
  properties, not PAC bounds or cryptographic proofs.
\item
  \textbf{High-impact coding-agent scenarios remain under-tested.} We
  have not yet fully benchmarked destructive git workflows and
  secret-exposure paths (e.g., accidental commit/exfiltration of
  repository credentials and CI tokens), where delayed detection may be
  especially costly.
\item
  \textbf{External validation gap.} Independent replication remains
  necessary for claims beyond controlled-environment demonstration.
\end{enumerate}

Benchmark-scope note: the aggregate metrics in Section 7 quantify
EDAMAME's downstream divergence/correlation performance under the
explicit-slice producer mode and do not directly measure
transcript-to-model accuracy of either the OpenClaw
\texttt{edamame-extrapolator} skill or EDAMAME's internal raw-session
ingest path. End-to-end reasoning-plane publication is validated
separately by the integration and E2E paths described in Section 7.7.

\subsubsection{9.2 Near-Term Work}\label{near-term-work}

\begin{itemize}
\tightlist
\item
  \textbf{Comparative evaluation on shared benchmarks}: run BadAgentUse
  scenarios against AgentSentinel, LlamaFirewall, and other monitors to
  produce head-to-head DSR/FPR comparisons. Requires adapting either the
  benchmark harness or the competing systems to a common telemetry
  interface.
\item
  \textbf{Compound multi-stage attacks}: scenarios where a single
  compromised action chain touches all five signal dimensions (network,
  Layer-7, posture, LAN, identity), validating cross-dimension
  divergence detection in a single evaluation.
\item
  \textbf{Adaptive evaluation}: extend BadAgentUse with scenarios that
  adapt their evasion strategy based on the monitor's known detection
  capabilities, following AgentSentinel's adaptive attack framework.
\item
  \textbf{Reasoning-plane auditing}: add a lightweight CoT consistency
  check as a complementary signal. Before accepting an intent
  declaration, evaluate whether it is consistent with the conversation
  context and previous declarations.
\item
  \textbf{Payload-aware detection}: extend Example Rule A to consider
  data volume and connection duration as additional signals for LotL
  detection.
\item
  \textbf{MCP-native poisoning scenarios}: extend BadAgentUse with
  scenario families inspired by MCPTox (tool-description poisoning) and
  MSB (tool-response injection). The two-plane approach should detect
  policy-inconsistent egress or file access induced by such attacks even
  when the MCP call itself appears legitimate, but this requires
  empirical validation.
\item
  \textbf{Observer-agnostic interface}: define an abstraction layer so
  that alternative telemetry providers (osquery, raw eBPF, Sysmon) can
  replace EDAMAME Posture, enabling independent replication and reducing
  the single-vendor dependency.
\item
  \textbf{Ablation baselines}: implement a telemetry-only baseline (B0:
  blacklist + anomaly score only, no intent gating) and an
  allowlist-only baseline (B1: no critical signals) to quantify the
  contribution of intent correlation beyond existing observer signals.
\item
  Measure rollback reliability under reversible action protocols.
\item
  Multi-host replication with federated observers.
\item
  Cross-platform benchmark runs on macOS and Windows (EDAMAME Posture
  supports both platforms; the benchmark harness requires adaptation).
\end{itemize}

\subsection{10. Ethics and Safe Use}\label{ethics-and-safe-use}

This system provides defensive runtime monitoring; it does not grant
authorization decisions or replace IAM/policy controls. It should be
deployed alongside baseline hardening controls (S1 and S2), manual
confirmation paths for high-impact actions, and claim discipline tied to
measured evidence and protocol scope.

\subsubsection{10.1 Telemetry Privacy and Data
Sovereignty}\label{telemetry-privacy-and-data-sovereignty}

EDAMAME Posture follows a \textbf{privacy-first, user-up} architecture:
all telemetry is collected, processed, and stored locally on the
endpoint. Session data, Layer-7 enrichment, and posture snapshots are
retained for a rolling 8-hour window and then discarded. No telemetry --
including process metadata, open file paths, command lines, or network
sessions -- is ever transmitted to the EDAMAME backend or any external
service. The EDAMAME platform (https://edamame.tech) is designed so that
endpoint security operates in a user/agent-up model where the user or
agent owns and controls their own security posture; no administrator can
look down into endpoints or collect personally identifiable information
(PII). This stands in contrast to traditional enterprise EDR/MDM
architectures where a central admin has full visibility into endpoint
activity.

In multi-user or shared-host deployments, Layer-7 enrichment may include
usernames and command-line arguments, creating a local privacy surface.
The MCP endpoint defaults to localhost binding with PSK authentication
(Section 10.2), ensuring telemetry is not exposed beyond the local host.

\subsubsection{10.2 MCP Endpoint Exposure}\label{mcp-endpoint-exposure}

The Streamable HTTP MCP endpoint defaults to localhost binding
(127.0.0.1:3000). Misconfiguration -- binding to 0.0.0.0 or routing
through a reverse proxy without authentication -- would expose sensitive
runtime telemetry to the network. EDAMAME Posture supports pre-shared
key (PSK) authentication; we recommend enabling it whenever the endpoint
is reachable beyond loopback. The observer process itself is part of the
attack surface (Section 4.3): an adversary who compromises the MCP
endpoint can suppress alerts or inject false telemetry. Deployment
hardening should follow the MCP specification's transport security
guidelines.

\subsubsection{10.3 Ethical Considerations (Menlo Report
Alignment)}\label{ethical-considerations-menlo-report-alignment}

We map our research to the Menlo Report's four ICT research ethics
principles:

\begin{itemize}
\tightlist
\item
  \textbf{Respect for Persons.} No human subjects or real-user data are
  involved. All evaluation runs in an isolated Lima VM with synthetic
  injections. No PII is collected, stored, or transmitted. The 8-hour
  local retention window (Section 10.1) ensures no persistent personal
  data accumulates.
\item
  \textbf{Beneficence.} The system is purely defensive. We document
  blind spots (LotL evasion, Section 4.3) and design limitations to
  prevent false confidence.
\item
  \textbf{Justice.} Benchmark scenarios, decision rules, and evaluation
  scripts are open-source for independent replication and critique.
\item
  \textbf{Respect for Law and Public Interest.} Telemetry collection
  occurs only on the researcher's own VM. No third-party systems are
  probed or scanned. Benchmark injections target localhost or controlled
  endpoints only.
\end{itemize}

\subsubsection{10.4 AI Tools Disclosure}\label{ai-tools-disclosure}

Generative AI tools (Claude, Cursor AI) assisted with code scaffolding
and paper editing. All outputs were reviewed and validated by the human
authors. No AI system is listed as an author. Scenario definitions,
decision rules, and architecture diagrams are human-authored. LLM-driven
test tiers (Categories A--C) use \texttt{openai/gpt-5.1} (with
compatible \texttt{azure-openai-responses/*} alternate provisioning);
deterministic tiers do not use LLM inference. Verdicts are deterministic
given fixed telemetry (Example Rule A); only the natural-language
explanation varies across runs.

\subsection{11. Conclusion}\label{conclusion}

We present a practical two-plane runtime security model: correlate agent
intent (reasoning plane) with host behavior (system plane), and flag
divergence when behavior is unexplained by intent or triggers
intent-independent critical signals (CVE vulnerability detection,
blacklisted traffic). We instantiate the model in OpenClaw with EDAMAME
Posture and evaluate it using trace-backed live artifacts.

The evidence supports the architectural claim: process-attributed
telemetry plus intent correlation yields operationally useful detection.
The boundary conditions are explicit: living-off-the-land behavior and
telemetry timing windows remain primary blind spots, and comparative
baseline studies are future work.

\subsection{Appendix A. Multi-Signal Scenario
Sketches}\label{appendix-a.-multi-signal-scenario-sketches}

These sketches illustrate how the same two-plane logic extends beyond
pure egress checks when richer telemetry is available.

\textbf{A1. Lateral movement (LAN + traffic).} A compromised skill
performs local recon and attempts service probing. LAN discovery APIs
detect new hosts/open ports while network sessions attribute traffic to
the originating process tree. If reasoning-plane intent remains
``document summarization,'' the mismatch yields divergence.

\textbf{A2. Credential theft (identity + traffic).} A supply-chain
compromise accesses credential stores and initiates outbound transfer.
Identity/breach signals increase risk context while Layer-7
open-file/process metadata and network telemetry provide direct evidence
for correlation.

\textbf{A3. Compound attack (network + Layer-7 + identity + posture).} A
staged workflow scans, reads sensitive material, and exfiltrates. No
single signal family is sufficient, but cross-signal correlation yields
stage-specific verdicts and supports clearer operator triage.

\subsection{Appendix B. Portability Notes and Future
Work}\label{appendix-b.-portability-notes-and-future-work}

Section 3.4 and Section 3.5 present the model-level portability argument
in the main narrative. This appendix keeps that discussion intentionally
brief: it documents the OpenClaw-specific correspondence and then scopes
portability work that remains future research.

\subsubsection{B.1 How the Skill Architecture Differs from Orchestration
Frameworks}\label{b.1-how-the-skill-architecture-differs-from-orchestration-frameworks}

Most agentic stacks are built on orchestration frameworks: LangGraph
(LangChain, 2024), CrewAI (Moura, 2024), and AutoGen (Microsoft, 2023).
These frameworks share several characteristics:

\begin{itemize}
\tightlist
\item
  \textbf{Graph-based execution.} Workflows compile into
  DAG/state-machine steps.
\item
  \textbf{Cloud-first telemetry.} Traces often flow through external
  observability services (e.g., LangSmith, AgentOps, Arize Phoenix).
\item
  \textbf{Framework coupling.} Guardrails are usually callback/runtime
  specific.
\end{itemize}

The OpenClaw skill model is architecturally distinct on every axis:

{\def\LTcaptype{none} % do not increment counter
\begin{longtable}[]{@{}
  >{\raggedright\arraybackslash}p{(\linewidth - 4\tabcolsep) * \real{0.1786}}
  >{\raggedright\arraybackslash}p{(\linewidth - 4\tabcolsep) * \real{0.4286}}
  >{\raggedright\arraybackslash}p{(\linewidth - 4\tabcolsep) * \real{0.3929}}@{}}
\toprule\noalign{}
\begin{minipage}[b]{\linewidth}\raggedright
Property
\end{minipage} & \begin{minipage}[b]{\linewidth}\raggedright
Orchestration frameworks
\end{minipage} & \begin{minipage}[b]{\linewidth}\raggedright
OpenClaw skill model
\end{minipage} \\
\midrule\noalign{}
\endhead
\bottomrule\noalign{}
\endlastfoot
Execution model & Compiled graph / state machine & Conversational
Markdown instructions \\
Runtime state & Framework-managed checkpoints & Agent-controlled memory
files \\
Security boundary & Framework callbacks / middleware & Declarative
manifest (\texttt{clawhub.json}) \\
Telemetry path & Cloud API (LangSmith, AgentOps) & Local-only MCP
(localhost, PSK) \\
Agent coupling & Framework-specific (LangGraph nodes, CrewAI tasks) &
Protocol-level (MCP tools) \\
Dependency weight & Heavy (Python ecosystem, SDK, cloud auth) &
Zero-framework (Markdown + MCP HTTP) \\
\end{longtable}
}

In OpenClaw, the skill is Markdown (\texttt{SKILL.md}) interpreted at
runtime, without compiled workflow graphs or framework SDK dependencies.
This keeps the security contract independent of framework internals.

\subsubsection{B.2 Skill-to-Architecture
Correspondence}\label{b.2-skill-to-architecture-correspondence}

The OpenClaw-facing reasoning skill and the EDAMAME-internal loops
(divergence plus posture/remediation) interact with three EDAMAME
layers, but through different interfaces depending on whether the
component runs inside or outside the observer process:

\textbf{Telemetry layer.} In the reference benchmark deployment, EDAMAME
Posture (via flodbadd) provides continual, process-attributed sessions
with Layer-7 enrichment. On workstation-hosted Cursor paths, the same
internal runtime is hosted by the EDAMAME Security app. The internal
divergence and vulnerability loops access this data directly via Rust
method calls (\texttt{get\_current\_sessions},
\texttt{get\_anomalous\_sessions}, \texttt{get\_blacklisted\_sessions},
\texttt{get\_lanscan}) within the host runtime; no MCP round-trip is
involved. This is a deliberate security property: the telemetry path
from capture to verdict is entirely internal, unreachable by external
processes (Section 4.4). The extrapolator skill does not read EDAMAME
telemetry at all; it reads only OpenClaw session transcripts via
platform APIs (\texttt{sessions\_list}, \texttt{sessions\_history}) and
publishes its behavioral model to EDAMAME through the
\texttt{upsert\_behavioral\_model} MCP tool.

\textbf{Scoring layer.} EDAMAME Foundation computes composite scores and
compliance mapping (CIS, ISO 27001, SOC 2). The internal engine reads
score directly (\texttt{get\_score}) for posture-delta detection. The
optional posture facade can read the same data via MCP
(\texttt{get\_score}, \texttt{advisor\_get\_todos}) while reusing the
same native remediation pipeline.

\textbf{Action layer.} EDAMAME's agentic pipeline
(\texttt{agentic\_execute\_action} / \texttt{advisor\_undo\_action})
provides execute-and-rollback with confirmation gates. The internal
posture loop owns scheduled remediation, while the optional posture
facade can consume this layer via MCP for explicit invocations. Example
Rule B intercepts here to enforce intent consistency before actions
execute.

The key insight is architectural alignment with security isolation: the
OpenClaw-facing reasoning skill interacts with EDAMAME through
authenticated MCP endpoints, while divergence and posture/remediation
loops run inside the observer and access telemetry directly. This
separation means that compromising the reasoning plane cannot tamper
with the telemetry-to-verdict path, while MCP remains the stable,
portable interface for external consumers.

\subsubsection{B.3 Future Work: Extending Beyond OpenClaw-Class
Flows}\label{b.3-future-work-extending-beyond-openclaw-class-flows}

This paper evaluates a specific class of agent behavior: cooperating
single-node conversational workflows sharing state through memory files.
We do \textbf{not} claim implemented portability to graph-orchestrated,
multi-node, or multi-agent frameworks in this draft.

Future work priority areas are:

\begin{enumerate}
\def\labelenumi{\arabic{enumi}.}
\tightlist
\item
  \textbf{Node-scoped intent declarations.} Move from session-scoped
  intent to per-node intent so allowlists can change safely across
  workflow steps.
\item
  \textbf{Workflow-aware correlation.} Attribute observations to the
  correct node across transitions and retries, rather than only to the
  currently active step.
\item
  \textbf{Delegation trust chains.} Verify that delegated intents are
  constrained by parent intent and do not expand permission budgets
  across agent boundaries.
\item
  \textbf{Monitorable workflow design.} Enforce explicit node intent,
  minimal per-node tool scope, observable transitions, bounded
  delegation depth, and immutable compiled graphs.
\item
  \textbf{Observer-agnostic replication.} Standardize interfaces so
  alternative observers can provide equivalent system-plane telemetry
  for independent studies.
\end{enumerate}

\subsubsection{B.4 Scope Boundary for This
Draft}\label{b.4-scope-boundary-for-this-draft}

Accordingly, this manuscript does not claim:

\begin{itemize}
\tightlist
\item
  validated LangGraph/CrewAI/AutoGen portability results,
\item
  formal workflow-level safety proofs for multi-node execution, or
\item
  cloud-governance analysis for deployments where telemetry leaves
  endpoints.
\end{itemize}

These remain future-work items after a workflow-aware prototype exists
and can be evaluated with trace-backed scenarios.

\subsection{12. Author Disclosure and Competing
Interests}\label{author-disclosure-and-competing-interests}

All two authors have financial relationships with EDAMAME Technologies
(founder, board member respectively). The bounding argument is that the
evaluation target is the correlation method over published artifacts,
not proprietary internals, and that independent replication is
practically enabled by open code, trace-level evidence, and observer
replaceability.

Frank Lyonnet is the founder of EDAMAME Technologies, the company that
develops EDAMAME Posture. This paper evaluates EDAMAME Posture as the
system-plane observer -- a core architectural dependency. Antoine
Clerget is a board member of EDAMAME Technologies.

To bound this conflict: (a) the observer is treated as a black-box
telemetry provider, not the evaluation target; (b) decision rules,
scenarios, harness, and scripts are open-source; (c) every verdict is
trace-backed with on-disk artifacts; and (d) the core observer
repositories are open-source: \textbf{flodbadd} (session capture +
Layer-7 attribution), \textbf{edamame\_foundation} (scoring/threat
infrastructure), and \textbf{threatmodels} (endpoint threat-model
definitions). All are published under EDAMAME Technologies and can be
audited or replaced by alternative observers. We actively invite
independent replication.

\subsection{References}\label{references}

\subsubsection{Attack Surface and
Benchmarks}\label{attack-surface-and-benchmarks}

\begin{itemize}
\tightlist
\item
  Kai Greshake, Sahar Abdelnabi, Shailesh Mishra, Christoph Endres,
  Thorsten Holz, Mario Fritz. ``Not what you've signed up for:
  Compromising Real-World LLM-Integrated Applications with Indirect
  Prompt Injection.'' arXiv:2302.12173 (2023).
  https://arxiv.org/abs/2302.12173
\item
  Qiusi Zhan, Zhixiang Liang, Zifan Ying, Daniel Kang. ``InjecAgent:
  Benchmarking Indirect Prompt Injections in Tool-Integrated Large
  Language Model Agents.'' arXiv:2403.02691 (2024).
  https://arxiv.org/abs/2403.02691
\item
  ``Exfiltration of personal information from ChatGPT via prompt
  injection.'' arXiv:2406.00199 (2024). https://arxiv.org/abs/2406.00199
\item
  Saeid Jamshidi et al.~``Securing the Model Context Protocol: Defending
  LLMs Against Tool Poisoning and Adversarial Attacks.''
  arXiv:2512.06556 (2025). https://arxiv.org/abs/2512.06556
\item
  Luo et al.~``MCPTox: A Benchmark for Tool Poisoning Attack on
  Real-World MCP Servers.'' arXiv:2508.14925 (2025).
  https://arxiv.org/abs/2508.14925
\item
  Wang et al.~``MCP Security Bench (MSB): Benchmarking Attacks Against
  Model Context Protocol in LLM Agents.'' arXiv:2510.15994 (2025).
  https://arxiv.org/abs/2510.15994
\item
  Zhou et al.~``Technical Report: Evaluating Goal Drift in Language
  Model Agents.'' arXiv:2505.02709 (2025).
  https://arxiv.org/abs/2505.02709
\end{itemize}

\subsubsection{Runtime Defense Systems}\label{runtime-defense-systems}

\begin{itemize}
\tightlist
\item
  ``AgentArmor: Enforcing Program Analysis on Agent Runtime Trace to
  Defend Against Prompt Injection.'' arXiv:2508.01249 (2025).
  https://arxiv.org/abs/2508.01249
\item
  ``LlamaFirewall: An open source guardrail system for building secure
  AI agents.'' arXiv:2505.03574 (2025). https://arxiv.org/abs/2505.03574
\item
  Li et al.~``AegisLLM: Scaling Agentic Systems for Self-Reflective
  Defense in LLM Security.'' arXiv:2504.20965 (2025).
  https://arxiv.org/abs/2504.20965
\item
  Wallarm et al.~``A2AS: Agentic AI Runtime Security and Self-Defense.''
  arXiv:2510.13825 (2025). https://arxiv.org/abs/2510.13825
\item
  ``MI9: An Integrated Runtime Governance Framework for Agentic AI.''
  arXiv:2508.03858 (2025). https://arxiv.org/abs/2508.03858
\item
  ``AgentSentinel: An End-to-End and Real-Time Security Defense
  Framework for Computer-Use Agents.'' arXiv:2509.07764 (2025).
  https://arxiv.org/abs/2509.07764
\item
  Zhang et al.~``Pro2Guard: Proactive Runtime Enforcement of LLM Agent
  Safety via Probabilistic Model Checking.'' arXiv:2508.00500 (2025).
  https://arxiv.org/abs/2508.00500
\item
  Chen et al.~``MCP-Guard: A Multi-Stage Defense-in-Depth Framework for
  Securing Model Context Protocol in Agentic AI.'' arXiv:2508.10991
  (2025). https://arxiv.org/abs/2508.10991
\item
  ``Securing AI Agent Execution.'' arXiv:2510.21236 (2025).
  https://arxiv.org/abs/2510.21236
\end{itemize}

\subsubsection{Formal Verification and Provable
Safety}\label{formal-verification-and-provable-safety}

\begin{itemize}
\tightlist
\item
  Xu et al.~``ShieldAgent: Shielding Agents via Verifiable Safety Policy
  Reasoning.'' arXiv:2503.22738 (2025). https://arxiv.org/abs/2503.22738
\item
  Li et al.~``VET Your Agent: Towards Host-Independent Autonomy via
  Verifiable Execution Traces.'' arXiv:2512.15892 (2025).
  https://arxiv.org/abs/2512.15892
\item
  Adharsh Kamath, Sishen Zhang, Calvin Xu, Shubham Ugare, Gagandeep
  Singh, Sasa Misailovic. ``Enforcing Temporal Constraints for LLM
  Agents (Agent-C).'' arXiv:2512.23738 (2025).
  https://arxiv.org/abs/2512.23738
\item
  Roham Koohestani. ``AgentGuard: Runtime Verification of AI Agents.''
  arXiv:2509.23864 (2025). https://arxiv.org/abs/2509.23864
\end{itemize}

\subsubsection{Chain-of-Thought
Monitoring}\label{chain-of-thought-monitoring}

\begin{itemize}
\tightlist
\item
  Rogers et al.~``Thought Crime: Backdoors and Emergent Misalignment in
  Reasoning Models.'' arXiv:2506.13206 (2025).
  https://arxiv.org/abs/2506.13206
\item
  Baker et al.~``Chain of Thought Monitorability: A New and Fragile
  Opportunity for AI Safety.'' arXiv:2507.11473 (2025).
  https://arxiv.org/abs/2507.11473
\item
  ``When Chain of Thought is Necessary, Language Models Struggle to
  Evade Monitors.'' arXiv:2507.05246 (2025).
  https://arxiv.org/abs/2507.05246
\item
  Goldowsky-Dill et al.~``CoT Red-Handed: Stress Testing
  Chain-of-Thought Monitoring.'' arXiv:2505.23575 (2025).
  https://arxiv.org/abs/2505.23575
\end{itemize}

\subsubsection{Agent Orchestration
Frameworks}\label{agent-orchestration-frameworks}

\begin{itemize}
\tightlist
\item
  LangChain / LangGraph. ``LangGraph: Build stateful, multi-actor
  applications with LLMs.'' (2024).
  https://langchain-ai.github.io/langgraph/
\item
  Joao Moura. ``CrewAI: Framework for orchestrating role-playing,
  autonomous AI agents.'' (2024). https://github.com/crewAIInc/crewAI
\item
  Qingyun Wu, Gagan Bansal, Jieyu Zhang, Yiran Wu, Beibin Li, Erkang
  Zhu, Li Jiang, Xiaoyun Zhang, Chi Wang. ``AutoGen: Enabling Next-Gen
  LLM Applications via Multi-Agent Conversation.'' arXiv:2308.08155
  (2023). https://arxiv.org/abs/2308.08155
\end{itemize}

\subsubsection{Standards and
Specifications}\label{standards-and-specifications}

\begin{itemize}
\tightlist
\item
  Edwin B. Wilson. ``Probable Inference, the Law of Succession, and
  Statistical Inference.'' JASA (1927).
  \href{https://www.tandfonline.com/doi/abs/10.1080/01621459.1927.10502953}{tandfonline.com/doi/10.1080/01621459.1927.10502953}
\item
  J. Sobel, D. Friedman. ``An Introduction to Reflection-Oriented
  Programming.'' (1996).
\item
  Model Context Protocol (MCP) Security Best Practices (rev 2025-06-18):
  \href{https://modelcontextprotocol.io/specification/2025-06-18/basic/security_best_practices}{modelcontextprotocol.io/\ldots/security\_best\_practices}
\item
  Model Context Protocol (MCP) Authorization (rev 2025-06-18):
  \href{https://modelcontextprotocol.io/specification/2025-06-18/basic/authorization}{modelcontextprotocol.io/\ldots/authorization}
\end{itemize}

\subsubsection{Incident Sources}\label{incident-sources}

\begin{itemize}
\tightlist
\item
  EDAMAME Posture (software):
  \href{https://github.com/edamametechnologies/edamame_posture}{github.com/edamametechnologies/edamame\_posture}
\item
  OpenClaw security advisories:
  \href{https://github.com/openclaw/openclaw/security/advisories}{github.com/openclaw/openclaw/security/advisories}
\item
  OpenClaw issue \#9627 (config secret sprawl):
  \href{https://github.com/openclaw/openclaw/issues/9627}{github.com/openclaw/openclaw/issues/9627}
\item
  CVE-2026-25253 (NVD):
  \href{https://nvd.nist.gov/vuln/detail/CVE-2026-25253}{nvd.nist.gov/vuln/detail/CVE-2026-25253}
\item
  CVE-2026-24763 (OpenCVE):
  \href{https://app.opencve.io/cve/CVE-2026-24763}{app.opencve.io/cve/CVE-2026-24763}
\item
  SecurityScorecard STRIKE exposure research (Moltbots):
  \href{https://securityscorecard.com/blog/beyond-the-hype-moltbots-real-risk-is-exposed-infrastructure-not-ai-superintelligence/}{securityscorecard.com/blog/\ldots{}}
\item
  declawed.io exposed instances dashboard:
  \href{https://declawed.io/}{declawed.io}
\item
  VirusTotal on weaponized agent skills (ClawHub):
  \href{https://blog.virustotal.com/2026/02/from-automation-to-infection-how.html}{blog.virustotal.com/2026/02/\ldots{}}
\item
  Palo Alto Networks Unit 42 (memory poisoning):
  \href{https://unit42.paloaltonetworks.com/indirect-prompt-injection-poisons-ai-longterm-memory/}{unit42.paloaltonetworks.com/\ldots{}}
\end{itemize}

\end{document}
